\documentclass[12pt]{article}
\usepackage{sbc-template}
\usepackage{graphicx,url}
\usepackage[utf8]{inputenc}
\usepackage[T1]{fontenc}      % glifos acentuados e hifenizacao correta no pdfLaTeX
\usepackage{lmodern}
\usepackage[brazil]{babel}
\usepackage{microtype}
\usepackage{booktabs}
\usepackage{textcomp}         % necessario para upquote nas listagens
\usepackage{listings}
\usepackage{xcolor}
\usepackage{amssymb}
\usepackage{amsmath}
\usepackage{stmaryrd}        % colchetes semanticos \llbracket \rrbracket
\sloppy

%% ── Configuração de listings ────────────────────────────────────────────────
\renewcommand{\lstlistingname}{Listagem}   % legenda em portugues
\definecolor{codegray}{rgb}{0.96,0.96,0.96}
\definecolor{commentcolor}{rgb}{0.40,0.40,0.40}
\definecolor{keywordcolor}{rgb}{0.13,0.29,0.53}
\definecolor{stringcolor}{rgb}{0.47,0.13,0.13}
\lstdefinestyle{haskell}{
  language=Haskell,
  backgroundcolor=\color{codegray},
  basicstyle=\ttfamily\footnotesize,
  keywordstyle=\color{keywordcolor}\bfseries,
  commentstyle=\color{commentcolor}\itshape,
  stringstyle=\color{stringcolor},
  breaklines=true,
  frame=single,
  numbers=left,
  numberstyle=\tiny\color{commentcolor},
  xleftmargin=1.6em,
  framexleftmargin=1.6em,
  captionpos=b,
  columns=flexible,            % evita espacamento artificial entre tokens
  keepspaces=true,
  upquote=true,                % crases e apostrofos retos (x `quot` y, modify')
  inputencoding=utf8,
  extendedchars=true,
  literate={á}{{\'a}}1 {ã}{{\~a}}1 {ç}{{\c c}}1 {é}{{\'e}}1 {í}{{\'i}}1
           {ó}{{\'o}}1 {õ}{{\~o}}1 {ú}{{\'u}}1 {â}{{\^a}}1 {ê}{{\^e}}1
}
\lstdefinestyle{ocaml}{
  language=[Objective]Caml,
  backgroundcolor=\color{codegray},
  basicstyle=\ttfamily\footnotesize,
  keywordstyle=\color{keywordcolor}\bfseries,
  commentstyle=\color{commentcolor}\itshape,
  stringstyle=\color{stringcolor},
  breaklines=true,
  frame=single,
  numbers=left,
  numberstyle=\tiny\color{commentcolor},
  xleftmargin=1.6em,
  framexleftmargin=1.6em,
  captionpos=b,
  columns=flexible,
  keepspaces=true,
  upquote=true,
  inputencoding=utf8,
  extendedchars=true,
  literate={á}{{\'a}}1 {ã}{{\~a}}1 {ç}{{\c c}}1 {é}{{\'e}}1 {í}{{\'i}}1
           {ó}{{\'o}}1 {õ}{{\~o}}1 {ú}{{\'u}}1 {â}{{\^a}}1 {ê}{{\^e}}1,
  morekeywords={effect,perform,match,continue,discontinue,with,fun,let,in,type,
                val,and,begin,end,try,raise,exception,function,module,open}
}
%% ────────────────────────────────────────────────────────────────────────────

\title{Mônadas e Efeitos Algébricos em Linguagens Funcionais:\\
Uma Análise Comparativa em Legibilidade, Composabilidade,\\
Poder de Abstração e Desempenho, Normalizada por\\
Linha de Base Intralinguagem}

\author{Diego Marchioni\inst{1}, Marco Rodrigo\inst{1}}

\address{Instituto de Ciências Exatas e Informática\\
  Pontifícia Universidade Católica de Minas Gerais (PUC Minas)\\
  Caixa Postal 1.686 -- 30.535-901 -- Belo Horizonte -- MG -- Brasil
  %% CONFERIR: dominio.
  \email{diego.marchioni@sga.pucminas.br, marco.rodrigo@pucminas.br}
}

\begin{document}
\maketitle

%% ============================================================
\begin{abstract}
Managing computational effects is a central challenge in functional language
design. Two approaches concentrate the attention of the literature: monads,
composed in Haskell via monad transformers, and algebraic effects with
handlers, supported natively in OCaml~5. Both structure impure computations but
differ in readability, composability, abstraction power, and runtime
performance. Since the effect mechanism cannot be isolated from the host
language runtime or syntax, the comparison is operationalized indirectly: each
mechanism is measured against an imperative baseline written in the same
language, and the two resulting ratios are then contrasted. The work combines a
literature-grounded conceptual analysis with four functionally equivalent
implementations of an expression interpreter with mutable state, exceptions and
output, a semantics-derived oracle test suite, a deterministic cross-language
workload generator, and an experimental protocol based on Criterion and
Bechamel. Counting lines of code over the implementations revealed that this
metric is highly sensitive to formatting style and to the explicit sequencing
OCaml requires, which led to restricting code-metric comparisons to
within-language pairs, mirroring the normalization already adopted for
performance. The artifacts exist as code and the empirical measurements are in
progress.
\end{abstract}

\noindent\textbf{Keywords:} algebraic effects; handlers; monads; monad
transformers; OCaml; Haskell; experimental evaluation.

\begin{resumo}
O gerenciamento de efeitos computacionais é um desafio central no design de
linguagens funcionais. Duas abordagens concentram a atenção da literatura: as
mônadas, compostas em Haskell por meio de \textit{monad transformers}, e os
efeitos algébricos com \textit{handlers}, com suporte nativo no OCaml~5. Ambas
estruturam computações impuras, mas diferem em legibilidade, composabilidade,
poder de abstração e desempenho. Como o mecanismo de efeitos não pode ser
isolado do \textit{runtime} nem da sintaxe da linguagem que o hospeda, a
comparação é operacionalizada de forma indireta: mede-se, para cada mecanismo,
o seu afastamento em relação a uma linha de base imperativa escrita na mesma
linguagem, e confrontam-se em seguida os dois afastamentos. O trabalho combina
uma análise conceitual fundamentada na literatura com quatro implementações
funcionalmente equivalentes de um interpretador de expressões com estado
mutável, exceções e saída, uma suíte de testes de oráculo derivada da semântica
formal, um gerador determinístico de cargas replicado nas duas linguagens e um
protocolo experimental baseado em Criterion e Bechamel. A contagem de linhas
sobre as implementações revelou que essa métrica é fortemente sensível ao
estilo de formatação e ao sequenciamento explícito exigido pelo OCaml, o que
levou a restringir as comparações de métricas de código a pares intralinguagem,
espelhando a normalização já adotada para o desempenho. Os artefatos existem
como código e as medições empíricas estão em curso.
\end{resumo}

\noindent\textbf{Palavras-chave:} efeitos algébricos; \textit{handlers};
mônadas; transformadores monádicos; OCaml; Haskell; avaliação experimental.

%% ============================================================
\section{Introdução}\label{sec:intro}
%% ============================================================

O paradigma funcional de programação tem como premissa central a transparência
referencial: a propriedade de que expressões com o mesmo valor são
intersubstituíveis em qualquer contexto sem alterar a denotação do programa
\cite{strachey:67}. Essa propriedade facilita o raciocínio equacional e a
verificação formal, e convive com os mecanismos de composição que sustentam a
modularidade do paradigma \cite{hughes:89}. Ela entra, porém, em tensão direta
com operações que um programa real precisa realizar: ler arquivos, lançar
exceções, manipular estado mutável. Tais operações são
denominadas \textit{efeitos computacionais} --- também referidos como
\textit{efeitos colaterais}, denominação que este trabalho evita por sugerir um
caráter acidental que a tradição algébrica rejeita --- e seu gerenciamento
constitui um dos problemas centrais do design de linguagens funcionais.

Para conciliar transparência referencial e expressividade prática, dois mecanismos
concentram a atenção da literatura e do design de linguagens contemporâneo. As
\textit{mônadas} foram introduzidas como estrutura semântica para uniformizar
noções de computação \cite{moggi:91}, e consolidadas como técnica de
programação que encapsula efeitos em tipos preservando a pureza do restante do
código \cite{wadler:92}. Em Haskell, são combinadas via \textit{monad
transformers} em pilhas como \texttt{StateT~s~(ExceptT~e~IO)} \cite{liang:95},
arranjo cuja ordenação estática no tipo é apontada como fonte de rigidez
estrutural \cite{kiselyov:13}. Os \textit{efeitos algébricos}, por sua vez,
declaram efeitos como conjuntos de operações \cite{plotkin:03} cuja semântica é
fixada por \textit{handlers} escopados \cite{plotkin:09}, o que permite
acrescentar ou substituir efeitos sem reestruturar código existente
\cite{pretnar:15}; objetivo semelhante é perseguido, por outra via, pelos
efeitos extensíveis sobre mônadas livres \cite{kiselyov:13}. Embora a teoria
seja definida algebricamente, a
implementação prática apoia-se em alguma forma de \emph{controle delimitado}:
tradução para continuações delimitadas em Koka \cite{leijen:17} e segmentos de
pilha alocados dinamicamente no OCaml~5 \cite{sivaramakrishnan:21}, distinção
retomada na Seção~\ref{sec:custo}.

\medskip
\noindent\textbf{Problema de pesquisa.} A relação teórica entre os dois
mecanismos está bem estabelecida, incluindo caracterizações formais de quando
efeitos algébricos modulares correspondem a transformadores monádicos
\cite{schrijvers:19,forster:17}. O mesmo não vale para as consequências
práticas dessa escolha. Um projetista que precise decidir entre as duas
abordagens não encontra na literatura, \emph{até onde alcançou a revisão
realizada}, evidência que trate \emph{simultaneamente e sobre um mesmo caso de
uso} as dimensões que governam a decisão: legibilidade, composabilidade, poder
de abstração e desempenho. Os estudos existentes concentram-se nos fundamentos
semânticos do mecanismo \cite{plotkin:03,plotkin:09}, na expressividade e na
composição de efeitos \cite{kiselyov:13,pretnar:15,lindley:17} ou no seu custo
de execução \cite{sivaramakrishnan:21,xie:20}, mas nenhum reúne as quatro
dimensões sobre um mesmo caso de uso; a Seção~\ref{sec:lacuna} detalha esse
levantamento e a Tabela~\ref{tab:lacuna} sintetiza a cobertura de cada
trabalho. O problema que este trabalho aborda é, portanto, a ausência de
evidência comparativa integrada e mensurável entre mônadas e efeitos
algébricos.

\medskip
\noindent\textbf{Justificativa.} A escolha entre os dois mecanismos deixou de
ser hipotética. Até recentemente, \textit{handlers} algébricos existiam
sobretudo em linguagens de pesquisa, e a decisão prática se reduzia a qual
arranjo monádico adotar; com a incorporação de \textit{handlers} nativos ao
OCaml~5 \cite{sivaramakrishnan:21}, a alternativa passou a estar disponível em
uma linguagem de uso industrial, com compilador otimizador e biblioteca madura.
Duas audiências enfrentam hoje essa decisão sem evidência que a informe. O
projetista de linguagens precisa saber o que se ganha e o que se perde ao
oferecer um dos mecanismos como construção primitiva. O engenheiro que já
mantém uma base monádica precisa estimar se a migração compensa, e a resposta
depende conjuntamente de quanto cada mecanismo custa em tempo de execução e de
quanto simplifica ou complica o código --- precisamente as dimensões que a
literatura trata em separado. Produzir essa evidência sobre um caso de uso
único, com critério de equivalência verificável e artefatos publicados, é o que
justifica o trabalho.

\medskip
\noindent\textbf{Objetivo geral.} Comparar mônadas e efeitos algébricos em
legibilidade, composabilidade, poder de abstração e desempenho, sobre um caso
de uso único que combine estado mutável, exceções e saída. O confundimento
entre mecanismo e linguagem hospedeira impede a medição direta, e a comparação
é feita por razões: cada mecanismo é medido contra uma \emph{linha de base
imperativa} na sua própria linguagem, e confrontam-se em seguida as duas razões
resultantes. Por linha de base imperativa entende-se uma implementação do
mesmo interpretador que realiza os três efeitos pelos meios diretos da
linguagem --- estado por célula mutável, erro por exceção nativa, saída por
escrita direta --- mantendo idênticas todas as demais escolhas, inclusive a
estrutura de dados do ambiente, de modo que a razão isole o custo do mecanismo.
O cotejo entre as duas linguagens é tratado como exploratório, pelas razões de
validade interna discutidas na Seção~\ref{sec:protocolo}.

\medskip
\noindent\textbf{Objetivos específicos.}
\begin{itemize}
  \item[\textbf{OE1}] Sistematizar as diferenças conceituais entre as duas
        abordagens e estabelecer o argumento de equivalência expressiva que
        torna a comparação justa.
  \item[\textbf{OE2}] Especificar formalmente a semântica do caso de uso e um
        critério verificável de equivalência funcional entre implementações.
  \item[\textbf{OE3}] Implementar quatro versões funcionalmente equivalentes do
        interpretador: mônadas em Haskell, efeitos algébricos em OCaml~5 e uma
        linha de base imperativa em cada linguagem.
  \item[\textbf{OE4}] Definir e aplicar métricas objetivas de legibilidade,
        composabilidade e poder de abstração sobre essas implementações, em
        pares intralinguagem.
  \item[\textbf{OE5}] Medir o sobrecusto de desempenho de cada mecanismo em
        relação à linha de base de sua própria linguagem, com tratamento
        explícito das ameaças à validade.
\end{itemize}
OE1 a OE3 são instrumentais: estabelecem a base conceitual, a referência formal
e os artefatos sobre os quais OE4 e OE5 operam. As questões de pesquisa
decorrem de OE4 (QP1--QP3) e de OE5 (QP4).

\medskip
\noindent\textbf{Questões de pesquisa.}
\begin{itemize}
  \item[\textbf{QP1}] Qual abordagem acrescenta menos carga estrutural de
        leitura sobre a linha de base da sua própria linguagem, medida pela
        extensão da lógica semântica, pelo número de operações de efeito e pela
        profundidade máxima de aninhamento (Seção~\ref{sec:metricas})? Caso os
        indicadores divirjam, a resposta é relatada por indicador, sem
        agregação.
  \item[\textbf{QP2}] Qual abordagem exige menos pontos de modificação para a
        composição e extensão de múltiplos efeitos, em relação à linha de base
        da sua própria linguagem (composabilidade)?
  \item[\textbf{QP3}] Qual abordagem acrescenta menos \textit{boilerplate} e
        obtém maior reuso de componentes sobre a linha de base da sua própria
        linguagem (poder de abstração)?
  \item[\textbf{QP4}] Qual é o custo em tempo de execução de cada abordagem em
        relação a uma linha de base imperativa \emph{na mesma linguagem}
        (desempenho)?
\end{itemize}
A legibilidade \emph{percebida} exigiria um estudo com participantes humanos e
está fora do escopo; a Seção~\ref{sec:metricas} explicita os \emph{proxies}
estruturais adotados em seu lugar. As hipóteses associadas são enunciadas na
Seção~\ref{sec:hipoteses} porque se formulam em termos das métricas definidas na
Seção~\ref{sec:metricas}.

\medskip
\noindent\textbf{Contribuições.} As contribuições deste trabalho são:
\begin{enumerate}
  \item Uma análise conceitual estruturada das diferenças entre mônadas e
        efeitos algébricos nas quatro dimensões propostas, incluindo o argumento
        de equivalência expressiva que sustenta a justiça da comparação
        (Seção~\ref{sec:fundamentacao}, em especial a
        Seção~\ref{sec:relacao}).
  \item Uma especificação semântica independente de implementação para o caso de
        uso, com critério de equivalência funcional verificável
        (Seção~\ref{sec:semantica}).
  \item Quatro implementações do interpretador, acompanhadas de uma suíte de
        testes de oráculo derivada da semântica que verifica a sua equivalência
        funcional, e de um gerador determinístico de cargas replicado nas duas
        linguagens (Seção~\ref{sec:impl}).
  \item Um desenho experimental que trata explicitamente o confundimento entre
        mecanismo e \textit{runtime}, normalizando cada mecanismo contra a linha
        de base de sua própria linguagem (Seção~\ref{sec:protocolo}).
\end{enumerate}

\medskip
\noindent\textbf{Estado do trabalho.} Este artigo dá continuidade ao relatório
de TCC~I \cite{marchioni:tcc1}. As contribuições 1, 2 e 4 estão
consolidadas; a contribuição 3 existe como código e está descrita na
Seção~\ref{sec:impl}, com a validação no ambiente definitivo em curso. As
medições empíricas ainda não foram realizadas: a Seção~\ref{sec:esperados}
apresenta os instrumentos de coleta e as observações já obtidas sobre o código,
e a Seção~\ref{sec:atividades} detalha as etapas restantes.

\medskip
O restante do artigo está organizado da seguinte forma. A
Seção~\ref{sec:fundamentacao} apresenta a fundamentação teórica. A
Seção~\ref{sec:relacionados} discute os trabalhos relacionados e delimita a
lacuna. A Seção~\ref{sec:metodologia} descreve a metodologia adotada. A
Seção~\ref{sec:impl} descreve as implementações desenvolvidas. A
Seção~\ref{sec:esperados} apresenta os instrumentos de coleta e as observações
preliminares. A Seção~\ref{sec:atividades} relata as atividades e o cronograma
restante. Por fim, a Seção~\ref{sec:consideracoes} apresenta as considerações
parciais.

\section{Fundamentação Teórica}\label{sec:fundamentacao}
%% ============================================================

Esta seção introduz os conceitos necessários para a comparação. Adota-se a
seguinte definição de trabalho: um \emph{efeito} é um conjunto de operações
cuja interpretação é deixada em aberto no ponto de uso e fixada em outro lugar
do programa. As duas abordagens comparadas diferem em \emph{onde} e \emph{como}
essa interpretação é fixada --- nas mônadas, estaticamente, pelo tipo da
computação e pela instância que o habita; nos efeitos algébricos,
dinamicamente, por um \textit{handler} instalado sobre um escopo. É essa
diferença que as quatro dimensões medem. A Seção~\ref{sec:fund-monadas} trata
das mônadas e de sua composição via transformadores; a
Seção~\ref{sec:fund-efeitos}, dos efeitos algébricos e dos \textit{handlers}; e
a Seção~\ref{sec:relacao} estabelece a relação formal de expressividade entre
as duas.

\subsection{Mônadas e Transformadores}\label{sec:fund-monadas}

As mônadas foram introduzidas como estrutura matemática para modelar noções de
computação de forma uniforme. Formalmente, dada uma categoria $\mathcal{C}$,
uma mônada é uma tripla $(T,\,\eta,\,\mu)$ em que $T:\mathcal{C}\to\mathcal{C}$
é um endofuntor e $\eta:\mathrm{Id}\Rightarrow T$ e $\mu:T^{2}\Rightarrow T$
são transformações naturais, sujeitas às leis das mônadas
($\mu \circ T\mu = \mu \circ \mu T$ e
$\mu \circ T\eta = \mu \circ \eta T = \mathrm{id}_{T}$). A contribuição central
está na separação entre \emph{valores}, de tipo $A$, e \emph{computações}, de
tipo $T A$, tornando o efeito acrescentado por $T$ passível de tratamento
uniforme no sistema de tipos \cite{moggi:91}.

As duas apresentações da mônada são equivalentes. A formulação em termos de
$(T,\eta,\mu)$ corresponde à \emph{tripla de Kleisli} $(T,\eta,(-)^{*})$, em
que a extensão $f^{*}: TA \to TB$ de $f: A \to TB$ é dada por
$f^{*} = \mu_{B} \circ Tf$. É essa segunda apresentação que o programador
manipula: $\eta$ é o \texttt{return} e $f^{*}$ é o \texttt{>>=} com os
argumentos trocados. As leis das mônadas traduzem-se exatamente nas três leis
que uma instância de \texttt{Monad} em Haskell deve satisfazer, e são elas que
garantem que o raciocínio equacional prometido na Seção~\ref{sec:intro}
sobreviva à introdução de efeitos.

A linhagem notacional das compreensões monádicas \cite{wadler:90} conduziu à
sintaxe \texttt{do} de Haskell, que é açúcar sintático sobre cadeias de
\texttt{>>=} \cite{wadler:92}. Cada efeito admite uma ou mais mônadas que o
modelam --- \texttt{State} para estado mutável, \texttt{Except} para exceções,
\texttt{IO} para entrada/saída, entre outras possíveis.

O problema surge ao combinar efeitos. Os \textit{monad transformers} foram
formalizados como mecanismo para construir mônadas compostas a partir de
mônadas mais simples, empilhando transformadores como \texttt{StateT} e
\texttt{ExceptT} sobre uma mônada base \cite{liang:95}. A composição, contudo,
impõe uma \emph{ordenação explícita e estática}, visível nos isomorfismos: com
o estado externo, $\texttt{StateT}\;s\;(\texttt{Except}\;e)\;a \cong s \to
(e + (a \times s))$, e o estado se perde no erro; com a exceção externa,
$\texttt{ExceptT}\;e\;(\texttt{State}\;s)\;a \cong s \to ((e + a) \times s)$, e
o estado sobrevive. Não é diferença estilística: são funções de tipos distintos.
A pilha \texttt{StateT~s~(ExceptT~e~IO)} é o exemplo canônico dessa dependência
de ordem e, precisamente por descartar o estado no erro, não é a adotada neste
trabalho (Seção~\ref{sec:impl-monadica}).

Cada camada exige ainda operações de elevação (\texttt{lift}, \texttt{liftIO})
para acessar efeitos de camadas internas, e a definição das instâncias de
elevação cresce quadraticamente com o número de transformadores e de classes.
Bibliotecas baseadas em classes de tipos (\textit{mtl}) eliminam grande parte
dessas elevações no ponto de uso, ao custo de deslocá-las para a definição das
instâncias; o custo estrutural da ordenação permanece intrínseco ao mecanismo.
Este trabalho adota \textit{mtl} deliberadamente, por comparar o transformador
\emph{como ele é idiomaticamente usado}; a escolha reduz de antemão a vantagem
esperada dos \textit{handlers} em volume de \textit{boilerplate} (H3), o que
torna mais significativa qualquer vantagem que ainda se observe.

\subsection{Efeitos Algébricos e \textit{Handlers}}\label{sec:fund-efeitos}

Os efeitos computacionais foram estabelecidos como operações algébricas sobre
uma teoria equacional, com a demonstração de que efeitos de interesse prático,
como estado, exceções, não determinismo e entrada/saída, podem ser
axiomatizados de forma uniforme \cite{plotkin:03}. Cada operação tem uma
assinatura e uma \emph{aridade}, isto é, o número de continuações que recebe:
$\texttt{get}$ tem aridade indexada pelos valores possíveis, $\texttt{set}$ tem
aridade um e $\texttt{raise}$ tem aridade \emph{zero} --- contrapartida
algébrica exata de o \textit{handler} de exceção descartar a continuação.

Nesse modelo, um efeito é declarado como um conjunto de operações, e a
composabilidade não é ausência de ordenação, e sim mudança de onde a ordenação
mora: ela deixa de ser fixada estaticamente no tipo da computação e passa a ser
determinada dinamicamente pelo escopo de instalação dos \textit{handlers}. Isso
a torna local e revisável sem alterar assinaturas, mas \textit{handlers} de
efeitos que interagem --- estado e exceção, por exemplo --- continuam não
comutando, do mesmo modo que as camadas de uma pilha monádica.

O modelo foi posteriormente estendido com \textit{handlers}, estruturas que
interceptam e interpretam operações de efeito dentro de um escopo delimitado
\cite{plotkin:09}. Um \textit{handler} recebe, a cada operação, a
\emph{continuação} da computação interrompida e decide se e como retomá-la.
Retomá-la uma vez modela estado e entrada/saída, descartá-la modela exceções, e
retomá-la múltiplas vezes modela não determinismo. É esse mecanismo de
continuação que confere aos \textit{handlers} seu poder expressivo. O modelo foi
consolidado na linguagem Eff \cite{bauer:15}, e, na linguagem Koka, um sistema
de \textit{effect rows} passou a verificar estaticamente quais efeitos uma
função pode realizar, garantindo que todos sejam tratados antes que a
computação seja considerada pura \cite{leijen:14}.

A ausência de um sistema de tipos para efeitos no OCaml~5, em que efeitos não
tratados só falham em tempo de execução, é identificada neste trabalho como a
principal desvantagem da abordagem adotada. Registre-se que a
\emph{segurança estática de efeitos} constitui, assim, uma quinta dimensão de
comparação: em Haskell, o tipo \texttt{ExceptT~String~(StateT~Env~IO)~Int}
declara os três efeitos e a omissão de tratamento é erro de compilação; em
OCaml~5, é erro de execução. Ela fica deliberadamente fora do arcabouço de
avaliação, que se limita às quatro dimensões enunciadas, e é retomada como
limitação na Seção~\ref{sec:protocolo}.

No OCaml~5, um efeito é introduzido por uma extensão do tipo \texttt{Effect.t},
cujo parâmetro registra o tipo de retorno da operação, e disparado por
\texttt{Effect.perform}. Um \textit{handler} é instalado com
\texttt{match\_with}, que recebe três cláusulas: \texttt{retc}, aplicada ao
valor quando a computação termina sem efeito; \texttt{exnc}, para exceções
nativas; e \texttt{effc}, a interpretação das operações, que devolve
\texttt{Some} para as operações que trata e \texttt{None} para as que repassa
ao \textit{handler} externo --- este último é o mecanismo de encadeamento entre
\textit{handlers}. A continuação \texttt{k} é manipulada com \texttt{continue}
(retoma), \texttt{discontinue} (retoma lançando exceção) ou simplesmente
descartada. Utilizam-se aqui \textit{handlers} \emph{profundos}
(\texttt{match\_with}), que permanecem instalados sobre a continuação retomada,
e não \emph{rasos} (\texttt{Effect.Shallow}), que precisariam ser reinstalados
a cada operação.

As continuações do OCaml~5 são \emph{one-shot}: cada continuação pode ser
retomada no máximo uma vez, restrição que é a contrapartida da implementação
por segmentos de pilha (Seção~\ref{sec:custo}). Efeitos que exigem retomada
múltipla, como o não determinismo, não são diretamente expressáveis nessa
implementação, restrição retomada na discussão de validade externa
(Seção~\ref{sec:protocolo}).

\subsection{Relação Formal entre as Abordagens}\label{sec:relacao}

A justiça da comparação proposta apoia-se em um argumento em quatro passos.

\emph{Primeiro}, a inter-expressividade entre \textit{handlers} de efeito,
reflexão monádica e controle delimitado foi estabelecida diretamente, com
traduções formais entre os três mecanismos e ressalvas explícitas para o
cenário tipado \cite{forster:17}. A noção relevante é a de
\emph{macro-expressividade} \cite{felleisen:91}: a pergunta não é se ambos
computam as mesmas funções, o que é trivialmente verdadeiro, mas se um pode ser
traduzido no outro por transformação local e composicional --- homomórfica nos
construtores sintáticos e conservativa sobre a linguagem-núcleo ---, sem
reescrita global do programa. O
resultado clássico de que qualquer mônada é representável por continuações
delimitadas e uma única célula de estado \cite{filinski:94} é o antecedente
histórico desse desenvolvimento.

\emph{Segundo}, a correspondência foi caracterizada de forma mais fina: há uma
classe identificável de efeitos algébricos \emph{modulares} em correspondência
com uma classe específica de transformadores monádicos \cite{schrijvers:19}.
Esse resultado delimita o que resta como diferença genuína: dentro dessa
classe, a escolha entre os mecanismos não altera o que pode ser expresso,
apenas como se expressa e a que custo. É exatamente esse resíduo, de natureza
ergonômica e de desempenho, que o presente trabalho mede.

\emph{Terceiro}, os três efeitos deste caso de uso são algébricos e usam a
continuação \emph{no máximo uma vez}. São algébricos porque cada um admite
apresentação por assinatura e equações com modelo livre:
\begin{itemize}
  \item \textbf{Estado}: operações \texttt{get} e \texttt{set}, com quatro
        equações --- duas leituras consecutivas devolvem o mesmo valor; ler após
        escrever devolve o escrito; a última escrita prevalece; escrever de volta
        o valor lido é inócuo. Modelo livre: $S \to (- \times S)$.
  \item \textbf{Exceções}: operação \texttt{raise} de aridade zero, sem
        equações; é a teoria livre sobre a assinatura. Modelo livre: $(- + E)$.
  \item \textbf{Saída}: operação \texttt{out}, cuja teoria é o monoide livre
        sobre os valores emitidos. Modelo livre: $(- \times V^{*})$.
\end{itemize}
Usam a continuação no máximo uma vez porque \texttt{get}, \texttt{set} e
\texttt{out} a retomam exatamente uma vez e \texttt{raise} a descarta --- uso
\emph{afim}, e não estritamente linear, justamente por causa do descarte. Nenhum
exige retomada múltipla, e todos cabem sob as continuações \emph{one-shot} do
OCaml~5. As equações do estado são, aliás, precisamente o que a suíte de oráculo
verifica nos casos de persistência e revinculação de \texttt{Let}
(Seção~\ref{sec:oraculo}).

\emph{Quarto}, segue-se que, \emph{para os três efeitos deste caso de uso}, os
dois mecanismos têm o mesmo alcance expressivo, e as diferenças observáveis são
de ergonomia (legibilidade, composabilidade, abstração) e de desempenho. A
ressalva é necessária: sob continuações \emph{one-shot}, efeitos que exigem
retomada múltipla não são diretamente expressáveis, de modo que a equivalência
aqui invocada é local ao caso de uso, e não uma equivalência entre os mecanismos
em geral.

Cabe ainda qualificar o alcance da correspondência. Ela recai sobre efeitos cuja
teoria admite apresentação \emph{algébrica} (de modelo livre) \cite{plotkin:03};
nem toda mônada é algébrica, sendo a mônada de continuação o contraexemplo
clássico, e os efeitos \emph{escopados} o caso limítrofe (Seção~\ref{sec:custo}).
É por essa razão --- e não por conveniência --- que o caso de uso adota ambiente
global sem restauração de vínculos (Seção~\ref{sec:caso-de-uso}).

%% ============================================================
\section{Trabalhos Relacionados}\label{sec:relacionados}
%% ============================================================

Enquanto a Seção~\ref{sec:fundamentacao} apresentou os mecanismos, esta seção
posiciona o presente trabalho frente à literatura, organizando-a em torno de
quatro eixos: abordagens alternativas de composição de efeitos, estudos de
implementação e expressividade de efeitos algébricos, o custo de execução dos
\textit{handlers}, e a perspectiva comparativa que motiva a lacuna preenchida
aqui.

\subsection{Abordagens Alternativas de Composição}\label{sec:alternativas}

A técnica \textit{data types à la carte}, baseada em pontos fixos de somas
abertas (coprodutos) de funtores, permite compor tipos de dados de forma
modular, com aplicação direta à composição de DSLs e interpretadores
\cite{swierstra:08}. Embora voltada à estruturação de tipos, e não de efeitos de
execução, é a precursora conceitual direta das abordagens de efeitos extensíveis
em Haskell: estas substituem o coproduto de funtores de sintaxe por um coproduto
de assinaturas de \emph{efeito}.

Os \textit{extensible effects} foram propostos como alternativa aos
transformadores monádicos, com base em uniões abertas de funtores que evita a
rigidez da pilha ordenada e permite interpretar efeitos em qualquer ordem
\cite{kiselyov:13}, técnica posteriormente refinada sobre mônadas
\textit{freer} \cite{kiselyov:15}. Embora implementada em Haskell, a proposta
antecipa vários aspectos da semântica de efeitos algébricos com
\textit{handlers}. A existência desse e de outros arcabouços de efeitos em
Haskell (por exemplo, \textit{polysemy}, \textit{fused-effects} e
\textit{effectful}) suscita a pergunta de por que comparar efeitos no OCaml com
mônadas em Haskell, e não as duas abordagens na mesma linguagem. A escolha,
justificada na Seção~\ref{sec:protocolo}, é comparar cada mecanismo \emph{como
ele é efetivamente utilizado e suportado em produção}.

\subsection{Implementação e Expressividade de Efeitos Algébricos}\label{sec:expressividade}

Uma implementação prática de \textit{handlers} em Haskell, baseada em mônadas
livres e continuações delimitadas, demonstrou a viabilidade do modelo em uma
linguagem sem suporte nativo, com desempenho competitivo em relação a código de
referência \cite{kammar:13}. Comparações posteriores entre \textit{handlers} e
\textit{monad transformers} \emph{na mesma linguagem} refinam esse quadro e
fornecem a base empírica mais próxima das hipóteses de desempenho adotadas aqui
\cite{xie:20}.

Um problema estrutural com \textit{handlers} no tratamento de efeitos escopados,
casos em que a semântica de um efeito depende do contexto de chamada, motivou a
proposta de \textit{scoped effects} como extensão do modelo de \textit{handlers}
\cite{wu:14}; variantes escopadas do efeito de estado, como estado local
restaurado ao fim de um bloco, recaem nesse regime e ficam deliberadamente fora
do caso de uso adotado, que emprega um ambiente global
(Seção~\ref{sec:caso-de-uso}). A linguagem Frank, por sua vez, generaliza a
abstração funcional, dispensando um construto separado de \textit{handler}, e
introduz \emph{multihandlers}, capazes de interpretar comandos de várias fontes
simultaneamente; seus autores argumentam que padrões como um operador
\texttt{pipe} ficam substancialmente mais simples em Frank do que com
\textit{handlers} unários em bibliotecas Haskell \cite{lindley:17}.

Uma introdução pragmática a efeitos algébricos e \textit{handlers}, por meio de
exemplos progressivos, discute explicitamente as implicações de legibilidade e
composabilidade que constituem dois dos eixos de comparação adotados aqui
\cite{pretnar:15}. Diferentemente do presente trabalho, essa introdução não
conduz benchmarks de desempenho nem compara as abordagens com métricas objetivas
sobre um caso de uso concreto.

\subsection{O Custo de Execução dos \textit{Handlers}}\label{sec:custo}

Uma linha de pesquisa que se intensificou a partir de 2021, e cujas raízes
remontam às primeiras implementações de \textit{handlers}, dedica-se a reduzir o
sobrecusto de execução do mecanismo, e seus resultados são relevantes para a
interpretação das medições deste trabalho. Um compilador otimizador da
linguagem Eff para OCaml, dirigido por informação de tipos e efeitos, mostrou
que boa parte desse sobrecusto pode ser eliminada por transformações
fonte-a-fonte, aproximando o desempenho do de código escrito à mão
\cite{karachalias:21}. Mais recentemente, avaliou-se a compilação
\textit{just-in-time} por rastreamento aplicada especificamente a efeitos e
\textit{handlers}, sob o argumento de que o fluxo de controle dinâmico que os
\textit{handlers} introduzem é precisamente aquilo que compiladores de
rastreamento otimizam bem \cite{gaissert:25}.

Duas consequências metodológicas decorrem dessa literatura. A primeira é que o
sobrecusto de \textit{handlers} é uma propriedade da \emph{implementação}, não
do mecanismo: medições obtidas em um \textit{runtime} específico não se
transferem para outro. A segunda é que o OCaml~5 ocupa uma posição particular
nesse espectro, por implementar \textit{handlers} nativamente sobre segmentos de
pilha alocados dinamicamente (\textit{fibers}), em vez de por tradução para
continuações \cite{sivaramakrishnan:21,ocaml:manual}. Note-se que as duas
propriedades estão ligadas: é porque as continuações são \emph{one-shot} que o
segmento de pilha não precisa ser copiado, de modo que a restrição registrada na
Seção~\ref{sec:fund-efeitos} é o preço, e também a razão, do baixo sobrecusto.
Ambas reforçam a decisão, adotada na Seção~\ref{sec:protocolo}, de reportar
razões de sobrecusto contra uma linha de base na mesma linguagem, em lugar de
tempos absolutos comparados entre linguagens.

\subsection{Perspectiva Comparativa e Lacuna Identificada}\label{sec:lacuna}

Três trabalhos são os mais diretamente relacionados. O primeiro descreve as
decisões de design do suporte nativo a efeitos algébricos no OCaml~5, com
microbenchmarks que mostram desempenho comparável ao de chamadas de função
diretas para casos comuns, tornando os efeitos viáveis para código de produção
\cite{sivaramakrishnan:21}; ele não compara, porém, efeitos algébricos com
mônadas, nem avalia legibilidade, composabilidade ou abstração. O segundo
compara \textit{handlers} e \textit{monad transformers} em desempenho dentro de
uma mesma linguagem \cite{xie:20}, eliminando o confundimento com o
\textit{runtime}, mas sem tratar as demais dimensões nem um caso de uso
integrado. O terceiro é o antecedente direto do caso de uso aqui adotado: os
interpretadores modulares construídos com transformadores monádicos
\cite{liang:95}, único trabalho da literatura revisada que, como este, articula
múltiplos efeitos sobre um caso de uso integrado. Ele não mede desempenho nem
legibilidade, e não podia contrastar transformadores com efeitos algébricos,
que só viriam a ser formulados como mecanismo de linguagem quase uma década
depois \cite{plotkin:03}; é dele, porém, que este trabalho herda a escolha do
interpretador como caso de uso. A Tabela~\ref{tab:lacuna} sintetiza a cobertura
dos trabalhos
discutidos e torna visível a lacuna que o presente trabalho busca preencher: uma
comparação conjunta das quatro dimensões sobre um mesmo caso de uso concreto.

\begin{table}[ht]
  \centering
  \caption{Cobertura dos trabalhos relacionados frente às quatro dimensões e ao
           critério de caso de uso integrado. $\bullet$ indica tratamento
           explícito; $\circ$, tratamento parcial ou incidental.}
  \label{tab:lacuna}
  \small
  \begin{tabular}{lccccc}
    \toprule
    \textbf{Trabalho} & \textbf{Legib.} & \textbf{Compos.} & \textbf{Abstr.}
      & \textbf{Desemp.} & \textbf{Caso integrado} \\
    \midrule
    \cite{plotkin:03,plotkin:09} &  & $\circ$ &  &  &  \\
    \cite{liang:95}              & $\circ$ & $\bullet$ & $\circ$ &  & $\bullet$ \\
    \cite{swierstra:08}          &  & $\bullet$ & $\bullet$ &  &  \\
    \cite{kiselyov:13,kiselyov:15} & $\circ$ & $\bullet$ & $\circ$ & $\circ$ &  \\
    \cite{kammar:13}             &  & $\circ$ &  & $\bullet$ &  \\
    \cite{pretnar:15}            & $\bullet$ & $\bullet$ & $\circ$ &  & $\circ$ \\
    \cite{lindley:17}            & $\circ$ & $\bullet$ & $\bullet$ &  &  \\
    \cite{sivaramakrishnan:21}   &  &  &  & $\bullet$ &  \\
    \cite{xie:20}                &  & $\circ$ &  & $\bullet$ &  \\
    \cite{karachalias:21,gaissert:25} &  &  &  & $\bullet$ &  \\
    \midrule
    \textbf{Este trabalho} & $\bullet$ & $\bullet$ & $\bullet$ & $\bullet$ & $\bullet$ \\
    \bottomrule
  \end{tabular}
\end{table}

%% ============================================================
\section{Metodologia}\label{sec:metodologia}
%% ============================================================

Esta seção descreve os procedimentos efetivamente adotados no desenvolvimento e
na avaliação do trabalho. A Seção~\ref{sec:caso-de-uso} delimita o caso de uso e
justifica a disciplina de escopo escolhida. A Seção~\ref{sec:semantica} fixa uma
semântica operacional independente de implementação e dela deriva o critério de
equivalência funcional que todas as versões devem satisfazer. A
Seção~\ref{sec:hipoteses} enuncia as hipóteses, a Seção~\ref{sec:metricas}
define as métricas, e a Seção~\ref{sec:protocolo} descreve o protocolo
experimental e as ameaças à validade.

Os procedimentos que diferem do que fora planejado no TCC~I estão enumerados na
Seção~\ref{sec:atividades}. Dois deles afetam diretamente a metodologia e são
detalhados adiante: a ordem das camadas da pilha monádica foi invertida, para
que o estado sobreviva ao erro como exige o critério de equivalência
(Seção~\ref{sec:impl-monadica}); e as comparações de métricas de código passaram
a ser feitas entre pares da mesma linguagem, porque a contagem sobre o código
efetivo mostrou a métrica sensível ao estilo de formatação e ao idioma
disponível em cada linguagem, e não ao mecanismo de efeitos
(Seção~\ref{sec:refinamento}). A verificação da equivalência funcional, antes
descrita apenas como intenção, é hoje operacionalizada por uma suíte de testes
executável (Seção~\ref{sec:oraculo}).

\subsection{Caso de Uso}\label{sec:caso-de-uso}

O caso de uso é um interpretador para uma linguagem de expressões simples
(Tabela~\ref{tab:syntax}) com suporte a aritmética inteira, vinculação e
leitura de variáveis em um ambiente mutável, sequenciamento e tratamento de
erros em tempo de execução. Esse cenário combina três efeitos que a semântica
apresenta em conjuntos de regras disjuntos (Seção~\ref{sec:semantica}) --- estado
mutável para o ambiente de variáveis, exceções para sinalização de erros
semânticos e saída para exibição de resultados intermediários ---, ainda que
interajam no comportamento observável, já que a terminação por erro delimita o
traço de saída.

Adota-se deliberadamente um ambiente único e global, sem restauração de
vínculos ao término de \texttt{Let}: a construção vincula (ou revincula) a
variável para todo o restante da avaliação. Além de manter o efeito de estado
simples, essa escolha decorre da delimitação teórica da Seção~\ref{sec:relacao}:
disciplinas de escopo com restauração correspondem a variantes \emph{escopadas}
do efeito de estado, que saem da classe algébrica sobre a qual repousa o
argumento de equivalência, e ficam como trabalho futuro.

Registre-se ainda uma consequência da sintaxe adotada. A linguagem não possui
condicional nem construto de captura de erro, de modo que o efeito de exceção é
exercitado apenas no regime de \emph{abortagem total}, sem recuperação: dos três
modos de uso da continuação identificados na Seção~\ref{sec:fund-efeitos}, o
caso de uso exercita a retomada única e o descarte, e não a retomada múltipla.
Conclusões sobre efeitos com captura em escopo ficam, portanto, fora do alcance
deste estudo (Seção~\ref{sec:protocolo}).

\begin{table}[ht]
  \centering
  \caption{Sintaxe abstrata da linguagem de expressões utilizada como caso de uso.}
  \label{tab:syntax}
  \small
  \begin{tabular}{rlp{0.46\linewidth}}
    \toprule
    \textbf{Expressão} & \textbf{Notação} & \textbf{Semântica} \\
    \midrule
    Literal         & $\texttt{Num}\;n$        & Inteiro $n$ \\
    Variável        & $\texttt{Var}\;x$        & Leitura do ambiente \\
    Adição          & $\texttt{Add}\;e_1\;e_2$ & Soma de $e_1$ e $e_2$ \\
    Multiplicação   & $\texttt{Mul}\;e_1\;e_2$ & Produto de $e_1$ e $e_2$ \\
    Divisão         & $\texttt{Div}\;e_1\;e_2$ & Divisão inteira truncada, erro se divisor $=0$ \\
    Ligação         & $\texttt{Let}\;x\;e\;b$  & Avalia $e$, vincula (ou revincula) $x$ no
                                                 ambiente global, avalia $b$; o vínculo
                                                 persiste após $b$ \\[2pt]
    Sequência       & $\texttt{Seq}\;e_1\;e_2$ & Avalia $e_1$, descarta, avalia $e_2$ \\
    Impressão       & $\texttt{Print}\;e$      & Imprime o valor de $e$ \\
    \bottomrule
  \end{tabular}
\end{table}

Esse cenário é recorrente na literatura de linguagens funcionais
\cite{liang:95,pretnar:15}, o que situa o caso de uso em terreno familiar. A
opção por um único caso de uso, em vez de múltiplos exemplos isolados, reduz a
interferência de fatores externos às abstrações de efeitos, ao custo de tornar
$n = 1$ toda generalização --- limitação retomada na validade externa
(Seção~\ref{sec:protocolo}).

\subsection{Semântica Operacional e Critério de Equivalência}
\label{sec:semantica}

Esta subseção fixa a semântica da linguagem da Tabela~\ref{tab:syntax} de forma
independente da implementação, de modo que as quatro versões possam ser
comparadas contra uma mesma referência formal --- e não umas contra as outras,
o que deixaria passar despercebido qualquer erro comum a duas delas.

\paragraph{Domínio semântico.}
Sejam $\mathcal{E}$ o conjunto das expressões geradas pela sintaxe da
Tabela~\ref{tab:syntax}, $\mathit{Var}$ o conjunto enumerável de identificadores
de variáveis e $\mathbb{Z}$ os inteiros. Um \emph{estado} (ou \emph{ambiente}) é
uma função parcial finita $\sigma : \mathit{Var} \rightharpoonup \mathbb{Z}$, e
$\Sigma$ denota o conjunto de estados. Um \emph{traço de saída} é uma sequência
$\omega \in \mathbb{Z}^{*}$. Seja $\mathit{Msg}$ a imagem das duas funções
canônicas de mensagem, $\mathit{msg\_unbound} : \mathit{Var} \to \mathit{Msg}$ e
$\mathit{msg\_div\_zero} \in \mathit{Msg}$. Identifica-se o inteiro $n$ com a
expressão $\texttt{Num}\;n$, de modo que valores são expressões. O conjunto de
\emph{configurações} é
\[
\mathcal{C} \;=\; (\mathcal{E} \times \Sigma \times \mathbb{Z}^{*})
\;\cup\;
(\{\bot_{\mathit{err}}\} \times \mathit{Msg} \times \Sigma \times \mathbb{Z}^{*}),
\]
em que $\bot_{\mathit{err}}$ marca a terminação por exceção. As configurações do
segundo componente da união, ditas \emph{configurações de erro}, são terminais:
nenhuma regra de transição parte delas. O conjunto dos \emph{resultados}
$\mathcal{R} \subseteq \mathcal{C}$ é o das configurações terminais: sucessos
$\langle v,\, \sigma',\, \omega \rangle$ com $v \in \mathbb{Z}$ e erros
$\langle \bot_{\mathit{err}},\, m,\, \sigma',\, \omega \rangle$.

Note-se que $\mathbb{Z}$ é o domínio da especificação, ao passo que as
implementações operam sobre inteiros de máquina de largura fixa --- 63 bits em
OCaml e 64 bits em Haskell. A equivalência funcional é, portanto, enunciada
sobre o subconjunto de $\mathbb{Z}$ representável em ambas as linguagens, e as
cargas de \textit{benchmark} são construídas para não sair dele
(Seção~\ref{sec:cargas}).

\paragraph{Relação de transição.}
A semântica é \emph{small-step}, com relação
$\Rightarrow \;\subseteq\; \mathcal{C} \times \mathcal{C}$.
A Figura~\ref{fig:semantica} apresenta as regras de redução e a gramática dos
contextos de avaliação \cite{felleisen:92}, da qual decorrem as regras
congruentes e a propagação de configurações de erro. A ordem de avaliação é
fixada à esquerda pela forma dos contextos: em $\texttt{Add}\;E\;e$, o operando
direito só é reduzido depois que o esquerdo se torna valor, e a decomposição de
uma configuração em contexto e redex é única. Daí o determinismo de
$\Rightarrow$.

\begin{figure}[ht]
\centering
\small
\fbox{\parbox{0.94\linewidth}{%
\textbf{Contextos de avaliação}
\par\smallskip
$E ::= \Box \mid \texttt{Add}\;E\;e \mid \texttt{Add}\;v\;E
      \mid \texttt{Mul}\;E\;e \mid \texttt{Mul}\;v\;E
      \mid \texttt{Div}\;E\;e \mid \texttt{Div}\;v\;E$
\par\smallskip
\hspace{2.1em}$\mid \texttt{Let}\;x\;E\;b \mid \texttt{Seq}\;E\;e
      \mid \texttt{Print}\;E$
\par\medskip\hrule\par\medskip
\textbf{(Var)}\quad
$x \in \mathrm{dom}(\sigma) \implies
 \langle \texttt{Var}\;x,\, \sigma,\, \omega \rangle
 \Rightarrow \langle \sigma(x),\, \sigma,\, \omega \rangle$
\par\medskip
\textbf{(Var-err)}\quad
$x \notin \mathrm{dom}(\sigma) \implies
 \langle \texttt{Var}\;x,\, \sigma,\, \omega \rangle
 \Rightarrow \langle \bot_{\mathit{err}},\,
 \mathit{msg\_unbound}(x),\, \sigma,\, \omega \rangle$
\par\medskip
\textbf{(Add)}\quad
$\langle \texttt{Add}\;n_1\;n_2,\, \sigma,\, \omega \rangle
 \Rightarrow \langle n_1+n_2,\, \sigma,\, \omega \rangle$
\par\medskip
\textbf{(Div)}\quad
$n_2 \neq 0 \implies
 \langle \texttt{Div}\;n_1\;n_2,\, \sigma,\, \omega \rangle
 \Rightarrow \langle n_1 \div n_2,\, \sigma,\, \omega \rangle$
\par\medskip
\textbf{(Div-zero)}\quad
$\langle \texttt{Div}\;n_1\;0,\, \sigma,\, \omega \rangle
 \Rightarrow \langle \bot_{\mathit{err}},\,
 \mathit{msg\_div\_zero},\, \sigma,\, \omega \rangle$
\par\medskip
\textbf{(Let)}\quad
$\langle \texttt{Let}\;x\;n\;b,\, \sigma,\, \omega \rangle
 \Rightarrow \langle b,\, \sigma[x \mapsto n],\, \omega \rangle$
\par\medskip
\textbf{(Seq)}\quad
$\langle \texttt{Seq}\;n\;e_2,\, \sigma,\, \omega \rangle
 \Rightarrow \langle e_2,\, \sigma,\, \omega \rangle$
\par\medskip
\textbf{(Print)}\quad
$\langle \texttt{Print}\;n,\, \sigma,\, \omega \rangle
 \Rightarrow \langle n,\, \sigma,\, \omega \cdot n \rangle$
\par\medskip
\textbf{(Cong)}\quad
$\langle e,\sigma,\omega\rangle \Rightarrow \langle e',\sigma',\omega'\rangle
 \implies \langle E[e],\sigma,\omega\rangle \Rightarrow
 \langle E[e'],\sigma',\omega'\rangle$
\par\medskip
\textbf{(Err-prop)}\quad
$\langle e,\sigma,\omega\rangle \Rightarrow
 \langle \bot_{\mathit{err}},m,\sigma',\omega'\rangle
 \implies \langle E[e],\sigma,\omega\rangle \Rightarrow
 \langle \bot_{\mathit{err}},m,\sigma',\omega'\rangle$
}}
\caption{Contextos de avaliação e regras da semântica operacional
\textit{small-step}. $\Rightarrow^{*}$ denota o fecho reflexivo-transitivo de
$\Rightarrow$; $\omega \cdot n$, a concatenação à direita; e $\div$, a divisão
inteira com truncamento em direção a zero. Configurações de erro são terminais.}
\label{fig:semantica}
\end{figure}

\paragraph{Progresso, terminação e função denotacional induzida.}
Duas propriedades sustentam a boa definição da semântica. \emph{Progresso}: por
indução na estrutura de $e$, toda configuração que não pertence a $\mathcal{R}$
admite transição, de modo que não há estados travados. \emph{Terminação}: a cada
passo, o tamanho da expressão decresce ou um redex é consumido, e a linguagem
não possui laços nem recursão; logo $\Rightarrow$ é fortemente normalizante.
Progresso, terminação e determinismo dão a totalidade e a unicidade da função
semântica $\mathcal{D} : \mathcal{E} \times \Sigma \to \mathcal{R}$, definida
por $\mathcal{D}(e, \sigma) = (v, \sigma', \omega)$ quando
$\langle e, \sigma, \varepsilon \rangle \Rightarrow^{*}
 \langle v, \sigma', \omega \rangle$ com $v \in \mathbb{Z}$, e por
$\mathcal{D}(e, \sigma) = (\bot_{\mathit{err}}, m, \sigma', \omega)$ quando a
redução termina em configuração de erro com mensagem $m$.

\paragraph{Ortogonalidade dos efeitos.}
Dizem-se \emph{ortogonais} os efeitos $E_{1}, \ldots, E_{n}$ quando suas regras
de transição podem ser apresentadas em conjuntos disjuntos $R_{1}, \ldots, R_{n}$
tais que, para cada $i \neq j$, as regras de $R_{i}$ alteram apenas
\emph{observáveis} --- resultado, estado, traço --- que as regras de $R_{j}$
preservam. No presente caso, a regra \textbf{(Print)} altera apenas $\omega$, a
regra \textbf{(Let)} altera apenas $\sigma$ (e \textbf{(Var)} apenas o consulta),
e as regras de erro (\textbf{Var-err}, \textbf{Div-zero}) alteram apenas o
componente de resultado. Essa disjunção é o que justifica, semanticamente, a
expectativa de que a inserção de um quarto efeito (Seção~\ref{sec:hipoteses},
H2) deva exigir intervenção apenas nas regras desse novo efeito, e não nas
existentes, expectativa que diferencia as duas abordagens em estudo.

\paragraph{Equivalência funcional entre implementações.}
Uma implementação $I$ é \emph{funcionalmente equivalente} a $\mathcal{D}$ se,
para toda expressão $e$ e todo estado inicial $\sigma_{0}$, vale
$\llbracket I \rrbracket(e, \sigma_{0}) = \mathcal{D}(e, \sigma_{0})$,
em que $\llbracket I \rrbracket$ denota a função produzida por $I$ sobre o
mesmo domínio $\mathcal{R}$. Em particular, exige-se: (i) o mesmo inteiro
retornado nos casos de sucesso, (ii) o mesmo estado final $\sigma'$,
(iii) o mesmo traço de saída $\omega$ (mesma sequência, na mesma ordem) e
(iv) a mesma mensagem de erro nos casos de falha, isto é, igualdade em
$\mathit{Msg}$, o que torna o critério verificável por comparação de cadeias.
A verificação é operacionalizada por uma suíte de testes de oráculo derivada de
$\mathcal{D}$, cuja construção é descrita na Seção~\ref{sec:oraculo}.

\subsection{Hipóteses}\label{sec:hipoteses}

As hipóteses são enunciadas em termos intralinguagem, em coerência com o escopo
das comparações fixado na Seção~\ref{sec:metricas}. Os critérios de decisão são
declarados aqui, antes da coleta, de modo que a confrontação com as evidências
não dependa de limiares escolhidos depois de conhecidos os dados.

\begin{itemize}
  \item[\textbf{H1}] (legibilidade) Em OCaml, o avaliador com \textit{handlers}
        apresentará contagem de operadores de efeito e profundidade de
        aninhamento próximas às da referência imperativa, com diferença não
        superior a 10\%; em Haskell, a versão monádica excederá as da sua base
        em \texttt{IORef}. A formulação intralinguagem substitui a comparação
        direta entre linguagens adotada no TCC~I, pelas razões discutidas na
        Seção~\ref{sec:refinamento}.
  \item[\textbf{H2}] (composabilidade) Em OCaml, a inserção do quarto efeito
        exigirá número de pontos de modificação próximo ao da linha de base; em
        Haskell, a versão monádica exigirá mais pontos que a sua base, por
        requerer alteração do tipo da pilha e das restrições de classe
        associadas \cite{kiselyov:13}. Registre-se que o uso de \textit{mtl}
        dispensa a reintrodução de elevações para as operações canônicas, de
        modo que a diferença esperada é menor do que a formulação do TCC~I
        sugeria.
  \item[\textbf{H3}] (abstração) Em cada linguagem, o mecanismo de efeitos
        apresentará volume de \textit{boilerplate} superior ao da sua linha de
        base, e o \textit{handler} permitirá reuso de componentes que a base não
        permite \cite{pretnar:15}. A adoção de \textit{mtl}
        (Seção~\ref{sec:fund-monadas}) reduz de antemão o \textit{boilerplate}
        da versão monádica.
  \item[\textbf{H4}] (desempenho) Os efeitos nativos do OCaml incorrerão em
        sobrecusto inferior a $1{,}5\times$ sobre a referência imperativa
        \emph{na mesma linguagem} \cite{sivaramakrishnan:21}, enquanto a pilha
        de transformadores em Haskell incorrerá em sobrecusto superior a
        $2{,}0\times$ sobre a sua própria linha de base (\texttt{IORef}) no
        cenário \textsc{State} \cite{xie:20}. Cada metade é considerada
        corroborada se o intervalo de confiança de 95\% da razão não contiver o
        limiar.
  \item[\textbf{P1}] (previsão exploratória) Espera-se que a razão de sobrecusto
        em OCaml seja inferior à razão em Haskell. Por comparar duas razões
        obtidas em \textit{runtimes} distintos, esta previsão é tratada como
        exploratória e considerada informativa apenas se os intervalos de
        confiança não se sobrepuserem (Seção~\ref{sec:protocolo}).
\end{itemize}

\subsection{Dimensões e Métricas de Avaliação}\label{sec:metricas}

\paragraph{Escopo e unidade de contagem.} Todas as métricas desta subseção são
comparadas \emph{entre pares da mesma linguagem} (efeitos \textit{versus}
imperativa em OCaml; mônadas \textit{versus} \texttt{IORef} em Haskell), e não
diretamente entre Haskell e OCaml; a justificativa empírica está na
Seção~\ref{sec:refinamento}. Fixam-se ainda duas unidades disjuntas: as métricas
de legibilidade (QP1) são contadas \emph{apenas sobre a função} \texttt{eval},
que é a lógica semântica, e as de \textit{boilerplate} (QP3) sobre o
\emph{arquivo menos} \texttt{eval}, que é o custo estrutural imposto pelo
mecanismo fora da lógica. Assim, nenhuma linha é contada nas duas dimensões.

\paragraph{Legibilidade (QP1).} Adotam-se três métricas estruturais. A primeira
é o total de linhas de código efetivas (LOC) de \texttt{eval}, excluindo linhas
em branco e comentários. A segunda é o número de \emph{operações de efeito},
definido de modo a ser comparável entre as versões de uma mesma linguagem: conta
tanto a realização abstrata quanto a direta do efeito --- em OCaml, cada
\texttt{perform} na versão com \textit{handlers} e cada acesso a
\texttt{Hashtbl}, \texttt{raise} ou emissão na base; em Haskell, cada
\texttt{>>=}, \texttt{<-} em bloco \texttt{do}, \texttt{throwError},
\texttt{get}, \texttt{modify'} e operação sobre \texttt{IORef}. A contagem é
reportada também por construtor semântico (são oito), para que versões mais
longas não sejam favorecidas. A terceira é a profundidade máxima de aninhamento
léxico, medida sobre a árvore de sintaxe: incrementam a profundidade os nós de
ligação (\texttt{let~\dots~in}, \texttt{<-} em bloco \texttt{do}), os braços de
casamento de padrão e os blocos condicionais. O aninhamento é adotado como
indicador de esforço de compreensão seguindo a \textit{Cognitive Complexity}
\cite{campbell:18}, que atribui incremento a cada nível de aninhamento
precisamente por a complexidade ciclomática \cite{mccabe:76} ser insensível a
ele. Validação empírica posterior, sobre 427 trechos de código e cerca de
24.000 avaliações humanas, encontrou correlação positiva da métrica com o tempo
de compreensão e com a avaliação subjetiva de inteligibilidade, ainda que com
resultados mistos quanto à correção das tarefas de compreensão
\cite{munozbaron:20}.

Essas métricas são \emph{proxies} de legibilidade, não medidas diretas da
legibilidade percebida; adotam-se múltiplos \emph{proxies} complementares. Um
estudo de legibilidade percebida, com participantes humanos, fica registrado
como trabalho futuro. Para
remover variação de estilo, o código é normalizado antes da contagem por
\texttt{ormolu} e por \texttt{ocamlformat} com perfil \texttt{conventional}
\cite{ormolu:tool,ocamlformat:tool}; os arquivos de configuração e as versões
exatas são publicados no repositório, dado que perfis distintos produzem
contagens distintas para o mesmo código. A análise qualitativa que acompanha as
métricas é guiada por uma rubrica fixada \emph{a priori}, com três critérios ---
(i) explicitação do fluxo de efeitos, (ii) localidade do tratamento de efeitos e
(iii) volume de cerimônia sintática ---, cada um pontuado em escala ordinal de
três pontos com âncoras descritas no documento de protocolo. A rubrica é
aplicada \emph{antes} de conhecidas as contagens, para que não constitua
racionalização dos números.

\paragraph{Composabilidade (QP2).} Avaliada pelo número de pontos de modificação
necessários para inserir um quarto efeito no interpretador existente. Para evitar
que a métrica favoreça uma das abordagens por construção, ``ponto de
modificação'' é definido a priori e de forma neutra como qualquer declaração ou
definição preexistente cujo texto precise ser alterado para acomodar o novo
efeito, excluindo-se as adições puras, que são contadas à parte. O quarto efeito
é especificado antes da implementação como leitura de um fluxo de entrada,
$\texttt{Input} : \texttt{unit} \to \texttt{int option}$, cujo esgotamento
produz erro; a escolha é deliberada por \emph{interagir} com o efeito de
exceção, e não ser meramente ortogonal, o que evita que H2 se confirme
trivialmente. A inserção é medida nas quatro implementações --- inclusive nas
linhas de base, que dão escala aos números --- e realizada em ramo separado do
repositório, de modo que os \textit{benchmarks} continuem a ser executados sobre
a versão de três efeitos descrita na Seção~\ref{sec:impl}.

\paragraph{Poder de abstração (QP3).} Avaliado pelo volume de
\textit{boilerplate}, na unidade fixada acima, e pelo número de instâncias de
reutilização de componentes, definido como o emprego de um mesmo
\textit{handler} ou camada em mais de um \emph{ponto de composição} --- isto é,
mais de um ponto do programa em que o mecanismo é instalado ou instanciado ---
sem reescrita. Registre-se que, com um único caso de uso, essa contagem tem
baixo poder discriminante; a resposta a QP3 apoia-se, portanto, sobretudo na
análise qualitativa das estratégias de reinterpretação de efeito
(Seção~\ref{sec:oraculo}). O critério de classificação de cada linha como
\textit{boilerplate} é fixado em um documento de protocolo, publicado com o
código, e aplicado por dupla contagem independente do autor e do orientador,
reportando-se o grau de concordância e as divergências.

\subsection{Protocolo Experimental e Ameaças à Validade}\label{sec:protocolo}

\paragraph{Configuração.} Os \textit{harnesses} de \textit{benchmark} estão
implementados com Criterion~1.6 \cite{criterion:08} em Haskell
(GHC~9.6.4, \texttt{-O2}) e Bechamel \cite{bechamel:21} em OCaml~5.2.1 com
\textit{flambda} (\texttt{-O3}), em quatro cenários. Dentro de cada linguagem, o
mecanismo e a sua linha de base são compilados com \textit{flags} idênticas. O
uso de um \textit{switch} com \textit{flambda} é condição necessária: sem ele, a
\textit{flag} \texttt{-O3} é silenciosamente ignorada, e a configuração descrita
aqui não corresponderia à efetivamente executada.

O cenário \textsc{Arith} avalia 10.000 expressões aritméticas sem variáveis. O
cenário \textsc{State} envolve 500 variáveis distintas e 10.000 leituras; o
prefixo de vinculação, um \texttt{Let} aninhado de 500 níveis, exercita
deliberadamente o crescimento dinâmico do segmento de pilha na versão com
\textit{handlers}. O cenário \textsc{Exn} avalia 10.000 expressões com 20\% de
divisões por zero. O cenário \textsc{Output} avalia 10.000 expressões com
\texttt{Print}, empregando o \textit{handler} acumulador em OCaml e um
\texttt{IORef} de lista em Haskell, em lugar de escrita real: mede-se assim o
custo do \emph{mecanismo} para o terceiro efeito sem que ele seja dominado pelo
subsistema de E/S do sistema operacional.

A geração das cargas é feita uma única vez, fora da região medida. Na versão
Haskell, utilizam-se \texttt{Control.Monad.State.Strict} e
\texttt{Data.Map.Strict}; os resultados aritméticos são forçados nas duas
versões, de modo que a disciplina de estrictness seja idêntica no numerador e no
denominador da razão, e os resultados finais são forçados à forma normal
(\texttt{nf}/\texttt{nfIO}). Nas versões OCaml, o resultado é envolvido em
\texttt{Sys.opaque\_identity} nas duas implementações, para que a otimização com
\textit{flambda} não elimine trabalho de forma assimétrica. Em ambas as
linguagens, o ambiente é reiniciado \emph{a cada iteração}, dentro da função
medida e igualmente nas duas versões, de modo que o custo da reinicialização se
cancele na razão; reiniciar apenas por amostra faria as iterações seguintes de
cada lote operarem sobre um ambiente já povoado.

O tratamento de exceções é feito por expressão avaliada, e não por carga
inteira, de modo que um erro no cenário \textsc{Exn} aborte apenas a expressão
corrente. Consequentemente, o \textit{handler} de exceção é instalado uma vez
por expressão, e o cenário mede também o custo de instalação e descarte de
\textit{handler} --- propriedade do uso realista do mecanismo, e não artefato,
mas que precisa ser considerada na interpretação.

Cada configuração é medida em 30 amostras, e os tempos são reportados com
intervalos de confiança obtidos por \textit{bootstrap}, procedimento nativo de
ambas as ferramentas \cite{georges:07}. Reporta-se também o volume de memória
alocada (via \texttt{GHC.Stats} e \texttt{Gc.quick\_stat}), grandeza
determinística que permite distinguir diferenças estruturais de ruído de
medição, e o coeficiente de determinação da regressão do Criterion; corridas com
$R^{2} < 0{,}95$ são descartadas e repetidas. Os três cenários são executados
com três sementes distintas, para verificar a estabilidade das razões.

%% TODO: registrar o tempo de parede da campanha completa apos a primeira
%% execucao de A4, e inserir pausas de resfriamento entre configuracoes se
%% necessario.
As medições são conduzidas em máquina com Intel Core i5-12500H (arquitetura
híbrida, quatro núcleos de desempenho e oito de eficiência), 32~GB de RAM, sob
Windows~11, com o plano de energia fixado em alto desempenho, o
\textit{turbo} desabilitado por \texttt{powercfg} (\textsc{perfboostmode}
igual a zero, limites mínimo e máximo do processador igualados), afinidade
restrita aos núcleos de desempenho por máscara de processador, verificação em
tempo real do antivírus e tarefas de indexação suspensas, nenhuma outra
aplicação ativa, e iterações de aquecimento para estabilização térmica e de
\textit{cache}.

\paragraph{Validação da equivalência funcional.} A entrada em produção dos
\textit{benchmarks} é condicionada à passagem integral da suíte de oráculo
descrita na Seção~\ref{sec:oraculo}, de modo que as diferenças observadas em
desempenho reflitam o custo do mecanismo de efeitos, e não divergências de
comportamento entre as implementações. A condição é reavaliada após qualquer
alteração no código, inclusive a inserção do quarto efeito.

\paragraph{Validade interna.} A principal ameaça é o \emph{confundimento entre o
mecanismo de efeitos e o \textit{runtime} da linguagem}: a versão monádica
executa sobre o GHC e as versões com efeitos e imperativa sobre o OCaml, com
modelos de coleta de lixo e otimização distintos. Comparar diretamente os tempos
absolutos entre Haskell e OCaml seria, portanto, comparar duas variáveis ao mesmo
tempo. Para mitigar, cada abordagem é normalizada contra uma linha de base
imperativa \emph{na sua própria linguagem}, a saber, a referência em OCaml
(Seção~\ref{sec:impl-base}) e uma referência quase imperativa em Haskell
baseada em \texttt{IORef} contendo o mesmo \texttt{Map~String~Int} da versão
monádica (e não uma tabela de espalhamento), de modo que a razão de sobrecusto
isole o mecanismo de efeitos, e não a estrutura de dados. Pela mesma razão, a
camada \texttt{IO} da pilha monádica, ociosa nos cenários sem \texttt{Print},
está presente também na base em \texttt{IORef} e se cancela na razão.

Uma ressalva quantitativa é devida na comparação entre as duas razões (P1). As
bases têm complexidades distintas --- \texttt{Hashtbl} em OCaml, com custo
amortizado constante, e \texttt{Map} em Haskell, logarítmico e com alocação por
consulta. Se o sobrecusto do mecanismo for aproximadamente aditivo por operação,
a razão vale $1 + \Delta / T_{\mathrm{base}}$, de modo que uma base mais cara
\emph{deflaciona} a razão. Como a base Haskell é a mais cara, o viés opera
\emph{contra} H4 e P1, que predizem razão maior naquela linguagem; caso a
predição se confirme, o resultado é mais forte do que o valor nominal sugere.

\paragraph{Ameaças específicas do ambiente de medição.} Três limitações do
ambiente disponível merecem registro. A primeira é o escalonamento em
arquitetura híbrida: a fixação de afinidade aos núcleos de desempenho reduz, mas
não elimina, a interferência do escalonador do sistema operacional. A segunda é
térmica: as medições são feitas em computador portátil, e cargas sustentadas ao
longo de trinta amostras podem induzir limitação por temperatura, fenômeno que
nenhuma configuração de energia elimina. Ambas afetam sobretudo os tempos
absolutos; como os resultados são reportados na forma de razões contra uma linha
de base medida nas mesmas condições e na mesma sessão, e a ordem de execução é
alternada entre tratamento e linha de base, o efeito sobre as conclusões é
atenuado. A terceira é a plataforma: as medições são feitas em Windows, ao passo
que a literatura de referência mede em Linux \cite{sivaramakrishnan:21}, e o
caminho de alocação de \textit{fibers} pode diferir entre as duas.

\paragraph{Validade de construto.} LOC, contagem de operações de efeito e
aninhamento são \emph{proxies} de legibilidade, e a Seção~\ref{sec:refinamento}
documenta um caso concreto em que a LOC se mostrou frágil. Registre-se ainda
que, na versão com \textit{handlers}, o estado reside em uma \texttt{Hashtbl} e
o traço em uma lista, ambos \emph{externos} ao \textit{handler}
(Seção~\ref{sec:impl-efeitos}); a alternativa canônica, em que estado e traço
são propagados pelo resultado do \textit{handler}, tem perfil de custo distinto,
de modo que o sobrecusto medido é o do mecanismo \emph{nesta} realização, e não
o de um \textit{handler} de estado em passagem de parâmetro.

\paragraph{Validade externa.} Os resultados restringem-se a um único caso de uso
(interpretador com três efeitos), a versões específicas de compilador e
hardware, e ao regime de \emph{abortagem total} do efeito de exceção
(Seção~\ref{sec:caso-de-uso}). Padrões de efeito não cobertos --- concorrência,
não determinismo (inexpressável diretamente sob as continuações
\emph{one-shot}) e efeitos escopados --- podem comportar-se de forma diferente.
Tampouco se avalia a \emph{segurança estática de efeitos}
(Seção~\ref{sec:fund-efeitos}), dimensão em que as duas abordagens diferem
substancialmente e que fica fora do arcabouço adotado. A escolha de um caso de
uso canônico na literatura e a publicação do código, dos dados brutos e das
configurações mitigam parcialmente essas limitações e favorecem a replicação.

\paragraph{Justificativa da escolha de linguagens.} Comparam-se efeitos
algébricos no OCaml~5 com mônadas em Haskell por serem, respectivamente, as
abordagens idiomáticas e suportadas em produção em cada linguagem. Uma
comparação intralinguagem (por exemplo, \textit{extensible effects} em Haskell
\textit{versus} transformadores \cite{xie:20}) é complementar e indicada como
trabalho futuro. Incluí-la aqui confundiria a maturidade e o nível de otimização
das bibliotecas com as propriedades do próprio mecanismo.

%% ============================================================
\section{Descrição das Implementações}\label{sec:impl}
%% ============================================================

Esta seção descreve o que foi desenvolvido desde a conclusão do TCC~I: as
quatro implementações do interpretador, o núcleo comum de sintaxe e geração de
cargas, a suíte de testes de oráculo e os dois \textit{harnesses} de
\textit{benchmark}. Registram-se também as decisões de implementação tomadas no
percurso, várias das quais corrigem problemas de equivalência funcional
identificados na revisão do TCC~I.

\subsection{Organização e Ferramental}\label{sec:organizacao}

O código está organizado em dois projetos independentes, um por linguagem, com
um núcleo comum de sintaxe e geração de cargas em cada um. A
Tabela~\ref{tab:artefatos} lista os artefatos e suas dimensões atuais.

\begin{table}[ht]
  \centering
  \caption{Artefatos desenvolvidos. LOC efetivas (sem linhas em branco nem
           comentários), antes da normalização por formatadores.}
  \label{tab:artefatos}
  \small\setlength{\tabcolsep}{3pt}
  \begin{tabular}{llr}
    \toprule
    \textbf{Artefato} & \textbf{Papel} & \textbf{LOC} \\
    \midrule
    \texttt{ocaml/lib/expr.ml}                  & sintaxe e geração de cargas & 52 \\
    \texttt{ocaml/effects/interp\_effects.ml}   & efeitos algébricos e \textit{handlers} & 123 \\
    \texttt{ocaml/imperative/interp\_imp.ml}    & linha de base do OCaml & 50 \\
    \texttt{ocaml/test/test\_oracle.ml}         & suíte de oráculo & 94 \\
    \texttt{ocaml/bench/bench\_main.ml}         & \textit{harness} Bechamel & 65 \\
    \texttt{haskell/src/Expr.hs}                & sintaxe e geração de cargas & 72 \\
    \texttt{haskell/src/Interp/Monadic.hs}      & \textit{monad transformers} & 54 \\
    \texttt{haskell/src/Interp/IORefBase.hs}    & linha de base do Haskell & 63 \\
    \texttt{haskell/test/Spec.hs}               & suíte de oráculo & 113 \\
    \texttt{haskell/bench/Main.hs}              & \textit{harness} Criterion & 37 \\
    \texttt{shared/conformance.json}            & artefatos comuns às duas linguagens & --- \\
    \texttt{docs/protocolo-metricas.md}         & protocolo de contagem & --- \\
    \midrule
    \multicolumn{2}{l}{\textbf{Total (código)}} & \textbf{723} \\
    \bottomrule
  \end{tabular}
\end{table}

O ferramental foi fixado da seguinte forma. No lado OCaml, um \textit{switch}
opam com OCaml~5.2.1 e \textit{flambda}, construído com \texttt{dune}~3.7 e a
biblioteca Bechamel para os \textit{benchmarks}. No lado Haskell, GHC~9.6.4 e
\texttt{cabal}, com \texttt{containers}, \texttt{mtl} e Criterion~1.6. A versão
do OCaml é determinada pela plataforma de desenvolvimento: o suporte nativo ao
Windows amadureceu com o opam~2.2 e a série 5.2 do compilador, de modo que a
versão inicialmente planejada (5.1.1) foi substituída. Nenhuma das
implementações depende de bibliotecas específicas de plataforma. A normalização
de estilo prevista nas métricas usa \texttt{ocamlformat} e \texttt{ormolu}, com
perfil e versão fixados no repositório.

Três artefatos precisam ser idênticos nas duas linguagens: o gerador de cargas,
as mensagens de erro canônicas e a tabela de casos do oráculo. Como a replicação
manual não oferece garantia de conformidade, os três são definidos em um arquivo
neutro (\texttt{shared/conformance.json}), lido por ambos os projetos no caso das
mensagens e dos casos de teste, e comparado por \textit{diff} contra um
\textit{dump} produzido por cada linguagem no caso do gerador. Um teste de
conformidade falha se as réplicas divergirem, o que torna o critério (iv) de
equivalência funcional verificável também \emph{entre} as linguagens, e não
apenas dentro de cada uma.

\subsection{Geração Determinística de Cargas}\label{sec:cargas}

Um problema não previsto no TCC~I emergiu ao especificar os cenários de
\textit{benchmark}: se cada linguagem gerasse suas expressões com o gerador
pseudoaleatório de sua biblioteca padrão, Haskell e OCaml avaliariam árvores de
sintaxe distintas, e a carga de trabalho passaria a ser uma segunda variável
confundida junto com o \textit{runtime}. A solução adotada foi implementar um
gerador congruencial linear próprio, $x_{n+1} = (1103515245\,x_n + 12345)
\bmod 2^{30}$, replicado de forma idêntica nas duas linguagens, com a semente
fixada no código e verificada pelo teste de conformidade
(Seção~\ref{sec:organizacao}). O módulo $2^{30}$ mantém todos os valores
intermediários dentro do inteiro nativo das duas linguagens.

Duas precauções são necessárias no uso do gerador. A primeira diz respeito aos
bits extraídos: em geradores congruenciais com módulo potência de dois, o bit
$k$ tem período $2^{k+1}$, de modo que os bits baixos ciclam rapidamente ---
tomar $x \bmod 8$ para escolher entre os oito construtores produziria uma
sequência de período 8. As escolhas discretas usam, por isso, os bits altos,
$(x \gg 15) \bmod n$. A segunda diz respeito à faixa dos literais: como as
árvores têm profundidade~3 e podem conter até oito folhas multiplicadas entre
si, os literais são sorteados em $[0,100]$, de modo que o pior caso, $100^{8} =
10^{16}$, permaneça dentro do inteiro de 63 bits comum às duas linguagens
(Seção~\ref{sec:semantica}).

Os quatro cenários são construídos a partir desse gerador. \textsc{Arith} produz
10.000 expressões aritméticas de profundidade~3 sem variáveis; para que o
cenário não gere erros, os divisores são literais no intervalo $[1,9]$.
\textsc{State} produz um \texttt{Let} aninhado que vincula 500 variáveis,
seguido de 10.000 expressões que leem duas variáveis sorteadas. \textsc{Exn}
produz 10.000 expressões nas quais 20\% terminam em uma divisão por zero.
\textsc{Output} produz 10.000 expressões contendo \texttt{Print}.

\subsection{Versão com Efeitos Algébricos}\label{sec:impl-efeitos}

Cada efeito é declarado como uma operação estendendo \texttt{Effect.t} e
disparado por \texttt{perform} (Listagem~\ref{lst:ocaml-eval}). Três decisões
merecem registro.

Primeiro, a operação \texttt{Get} devolve \texttt{int option}, e não
\texttt{int}. A consequência é que a política de erro para variáveis não
vinculadas é decidida no avaliador, que dispara \texttt{Fail} com a mensagem
canônica, e o \textit{handler} de estado permanece ignorante da semântica de
erro. Na formulação inicial do TCC~I, \texttt{Get} devolvia \texttt{int} e a
busca na tabela lançava \texttt{Not\_found}, o que quebrava a equivalência
funcional.

Segundo, a ordem de avaliação dos operandos é fixada com \texttt{let}
explícitos. O OCaml não especifica a ordem de avaliação dos argumentos de uma
aplicação, de modo que \texttt{eval~a~+~eval~b} não garante que os efeitos de
\texttt{a} precedam os de \texttt{b}, e o critério (iii) de equivalência
funcional poderia falhar de forma dependente do compilador.

Terceiro, o efeito de exceção é modelado descartando a continuação, e o de
estado retomando-a (Listagem~\ref{lst:ocaml-handler}). Nesta implementação, os
dois \textit{handlers} podem ser aninhados em qualquer ordem, por duas razões
que convém não confundir com uma propriedade geral do mecanismo: o estado reside
em uma \texttt{Hashtbl} \emph{externa} ao \textit{handler}, e não no seu
resultado; e \texttt{run\_state} é polimórfico no tipo de resposta, ao passo que
\texttt{run\_exn} o fixa. \textit{Handlers} que transformam o resultado impõem
restrição de ordem tanto quanto uma camada monádica.

\begin{lstlisting}[style=ocaml,caption={Declaração dos efeitos e casos representativos do avaliador em OCaml~5.},label={lst:ocaml-eval}]
type _ Effect.t +=
  | Get    : string -> int option Effect.t
  | Set    : string * int -> unit Effect.t
  | Fail   : string -> 'a Effect.t
  | Output : int -> unit Effect.t

let rec eval : expr -> int = function
  | Var x -> (match perform (Get x) with
              | Some v -> v
              | None -> perform (Fail (msg_unbound x)))
  | Div (a, b) ->
      let x = eval a in              (* ordem fixada com let *)
      let y = eval b in
      if y = 0 then perform (Fail msg_div_zero) else x / y
  | Let (x, e, b) -> let v = eval e in perform (Set (x, v)); eval b
  (* Num, Add, Mul, Seq, Print analogos *)
\end{lstlisting}

\begin{lstlisting}[style=ocaml,caption={\textit{Handlers} de estado (retoma a continuação) e de exceção (descarta a continuação), e sua composição.},label={lst:ocaml-handler}]
let run_state (tbl : (string, int) Hashtbl.t) (main : unit -> 'a) : 'a =
  match_with main ()
    { retc = (fun v -> v)
    ; exnc = (fun e -> raise e)
    ; effc = (fun (type a) (eff : a Effect.t) ->
        match eff with
        | Get x      -> Some (fun (k : (a, _) continuation) ->
                           continue k (Hashtbl.find_opt tbl x))
        | Set (x, v) -> Some (fun (k : (a, _) continuation) ->
                           Hashtbl.replace tbl x v; continue k ())
        | _          -> None) }

let run_exn (main : unit -> int) : (int, string) result =
  match_with main ()
    { retc = (fun v -> Ok v)
    ; exnc = (fun e -> raise e)
    ; effc = (fun (type a) (eff : a Effect.t) ->
        match eff with
        | Fail msg -> Some (fun (_ : (a, _) continuation) -> Error msg)
        | _        -> None) }

(* Composicao usada na execucao e nos benchmarks *)
let run (tbl : (string, int) Hashtbl.t) (e : expr) : (int, string) result =
  run_state tbl (fun () -> run_output (fun () -> run_exn (fun () -> eval e)))
\end{lstlisting}

Um terceiro \textit{handler} interpreta \texttt{Output}. Foram implementadas
duas versões: uma que escreve em \texttt{stdout}, usada na execução normal, e
uma que acumula os valores emitidos em uma lista, usada pela suíte de oráculo e
pelo cenário \textsc{Output}. Essa é uma manifestação concreta da propriedade
que o trabalho se propõe a avaliar: trocar a interpretação de um efeito exigiu
apenas fornecer outro \textit{handler}, sem alterar uma linha do avaliador.

\subsection{Versão Monádica}\label{sec:impl-monadica}

Os três efeitos são combinados na pilha
\texttt{ExceptT~String~(StateT~Env~IO)}, com \texttt{Env~=~Map~String~Int}
(Listagem~\ref{lst:haskell}). A ordem das camadas é uma decisão semântica, e não
estilística: com \texttt{ExceptT} externo, o estado sobrevive ao lançamento da
exceção, como exige o critério (ii) de equivalência funcional e como se comporta
a versão com efeitos, cujo estado vive em uma tabela externa. A pilha inversa,
\texttt{StateT~Env~(ExceptT~String~IO)}, descartaria o estado no erro, e essa
divergência é justamente um dos casos cobertos pela suíte de oráculo
(Seção~\ref{sec:oraculo}).

A divisão usa \texttt{quot}, e não \texttt{div}: o primeiro trunca em direção a
zero, coincidindo com o operador \texttt{/} do OCaml e com a semântica fixada na
Seção~\ref{sec:semantica}, enquanto \texttt{div} arredonda para baixo e
divergiria com operandos negativos. Utilizam-se
\texttt{Control.Monad.State.Strict} e \texttt{Data.Map.Strict}; as atualizações
de estado empregam \texttt{modify'} e os resultados aritméticos são forçados
explicitamente, para que a disciplina de estrictness seja idêntica à da linha de
base e a medição não capture acúmulo de \textit{thunks} em lugar do custo do
mecanismo.

\begin{lstlisting}[style=haskell,caption={Interpretador em Haskell com \textit{monad transformers}.},label={lst:haskell}]
type Env = Map String Int
type Interp a = ExceptT String (StateT Env IO) a

eval :: Expr -> Interp Int
eval (Num n) = return n
eval (Var x) = do
  env <- get
  maybe (throwError (msgUnbound x)) return (Map.lookup x env)
eval (Add a b) = do
  x <- eval a
  y <- eval b
  return $! x + y
eval (Div a b) = do
  x <- eval a
  y <- eval b
  when (y == 0) (throwError msgDivZero)
  return $! x `quot` y
eval (Let x e b) = do
  v <- eval e
  modify' (Map.insert x v)
  eval b
eval (Seq a b) = eval a >> eval b
eval (Print e) = do
  v <- eval e
  emit v
  return v
\end{lstlisting}

\subsection{Linhas de Base}\label{sec:impl-base}

A referência imperativa em OCaml usa \texttt{Hashtbl} para o ambiente, exceções
nativas para erros e escrita direta para a saída, com as mesmas mensagens
canônicas e a mesma fixação explícita da ordem de avaliação da versão com
efeitos. A tabela é recebida como parâmetro, e não capturada de um escopo
global, para que o padrão de acesso seja idêntico ao da versão com
\textit{handlers}. Seu avaliador é estruturalmente idêntico ao da
Listagem~\ref{lst:ocaml-eval}: a diferença reduz-se a substituir
\texttt{perform~(Get~x)} por \texttt{Hashtbl.find\_opt~env~x},
\texttt{perform~(Set~(x,v))} por \texttt{Hashtbl.replace~env~x~v} e
\texttt{perform~(Fail~m)} por \texttt{raise~(Runtime\_error~m)}. É esse cotejo
que isola o que o mecanismo de efeitos acrescenta.

A linha de base do Haskell usa um \texttt{IORef} contendo o mesmo
\texttt{Map~String~Int} da versão monádica, com erros por exceção
(\texttt{throwIO}) e o mesmo forçamento explícito dos resultados aritméticos. A
escolha do \texttt{Map}, em lugar de uma tabela de espalhamento, é deliberada:
sendo a estrutura de dados idêntica à da versão monádica, a razão entre as duas
isola o custo do mecanismo de efeitos, e não o da representação do ambiente.

Em ambas as linguagens, a emissão de saída é feita através de uma referência a
função (\texttt{emit}) que a suíte de oráculo substitui temporariamente. A
uniformidade é deliberada: ela dispensa a captura de \texttt{stdout} na suíte
Haskell, que dependia do \textit{buffering} do descritor e observava a
\emph{renderização} do traço em lugar do traço $\omega$ definido na
Seção~\ref{sec:semantica}.

\subsection{Suíte de Testes de Oráculo}\label{sec:oraculo}

A suíte materializa o critério de equivalência funcional da
Seção~\ref{sec:semantica}. Cada caso fixa uma tripla esperada, composta por
resultado, ambiente final ordenado e traço de saída, e cada implementação é
verificada contra a mesma tabela de casos, definida no arquivo de conformidade e
lida pelas duas linguagens. São 25 casos por implementação, cobrindo todos os
construtores da Tabela~\ref{tab:syntax} em situações de sucesso e de erro,
incluindo estados iniciais não vazios e operandos negativos.

Quatro grupos de casos merecem destaque por travarem decisões de projeto, e não
apenas por verificarem comportamento. O caso da divisão truncada,
$\texttt{Div}(-7, 2)$, usa operandos de sinais opostos, única configuração em
que \texttt{quot} e \texttt{div} divergem ($-3$ contra $-4$), e portanto a única
capaz de detectar a troca. O caso de ordem de avaliação usa \texttt{Print} nos
dois operandos, única configuração que torna a ordem observável, e falha se o
compilador avaliar o operando direito primeiro. Os casos de persistência e
revinculação de \texttt{Let} falham se alguma implementação introduzir
restauração de escopo. O caso em que estado e saída sobrevivem ao erro falha se
a pilha monádica for invertida, o que transforma uma decisão sutil de ordenação
de camadas em uma verificação automática.

A observação do traço $\omega$ é feita sem modificar os avaliadores. Na versão
com efeitos, por um \textit{handler} alternativo de \texttt{Output}; nas demais,
substituindo a referência a função de emissão. A assimetria é, ela própria, um
dado qualitativo relevante para QP3: nas três versões sem \textit{handlers} foi
preciso \emph{inserir um ponto de indireção antecipando a necessidade}, ao passo
que na versão com \textit{handlers} nada no avaliador precisou mudar, porque a
abstração já previa a substituição do interpretador do efeito.

\subsection{\textit{Harnesses} de \textit{Benchmark}}\label{sec:harness}

Os dois \textit{harnesses} implementam os quatro cenários da
Seção~\ref{sec:protocolo}, com as precauções ali descritas: geração das cargas
fora da região medida, reinicialização do ambiente a cada iteração e igualmente
nas duas versões de cada linguagem, forçamento à forma normal em Haskell,
\texttt{Sys.opaque\_identity} em OCaml, e tratamento de exceções por expressão
avaliada.

\subsection{Estado Atual e Pendências}\label{sec:estado}

As quatro implementações, a suíte de oráculo e os dois \textit{harnesses}
existem como código. A validação sobre o ambiente definitivo está em andamento e
é pré-requisito para a execução dos \textit{benchmarks}, conforme o cronograma
da Seção~\ref{sec:atividades}. Três pendências técnicas estão identificadas: o
relatório do \textit{harness} Bechamel depende de detalhes de API que variam
entre versões da biblioteca, razão pela qual a versão é fixada no \texttt{opam};
as anotações de tipo dos \textit{handlers} polimórficos em OCaml podem exigir
ajuste fino; e o teste de conformidade entre as réplicas do gerador precisa ser
executado nas duas plataformas. Soma-se a dependência de que o \textit{switch}
com \textit{flambda} seja construído com êxito; caso não seja, a alternativa
registrada é medir sem \textit{flambda}, declarando explicitamente a mudança de
configuração. Nenhuma dessas pendências afeta o desenho experimental, apenas a
sua execução.

%% ============================================================
\section{Instrumentos de Coleta e Observações Preliminares}\label{sec:esperados}
%% ============================================================

As tabelas de coleta desta seção são \emph{modelos} a serem preenchidos após a
execução descrita na Seção~\ref{sec:atividades}. Nenhum valor de desempenho é
reportado neste ponto de controle. Reporta-se, contudo, uma observação obtida
sobre o código já escrito, que motivou um refinamento da metodologia.

\subsection{Refinamento das Métricas de Legibilidade}\label{sec:refinamento}

No TCC~I, a contagem preliminar de LOC foi feita sobre \emph{excertos
ilustrativos} das implementações e apontava a versão com efeitos como a mais
enxuta. Refeita a contagem sobre as implementações efetivas, o resultado se
inverte (Tabela~\ref{tab:prelim-loc}): a versão monádica em Haskell apresenta a
menor contagem, e a ordenação resultante não separa os mecanismos de efeito das
linhas de base.

%% TODO: preencher a coluna "Pos-normalizacao" apos executar ormolu e
%% ocamlformat sobre o codigo efetivo (atividade A3).
\begin{table}[ht]
  \centering
  \caption{LOC efetivas da função \texttt{eval} nas implementações
           desenvolvidas, antes e depois da normalização por formatadores.}
  \label{tab:prelim-loc}
  \small
  \begin{tabular}{llcc}
    \toprule
    \textbf{Versão} & \textbf{Linguagem} & \textbf{Pré-normalização}
      & \textbf{Pós-normalização} \\
    \midrule
    Efeitos algébricos & OCaml~5  & 30 & -- \\
    Imperativa         & OCaml    & 29 & -- \\
    Mônadas            & Haskell  & 21 & -- \\
    \texttt{IORef}     & Haskell  & 30 & -- \\
    \bottomrule
  \end{tabular}
\end{table}

Duas afirmações distintas devem ser separadas aqui. A primeira é que a contagem
do TCC~I estava errada, e a razão é conhecida: ela mediu excertos, e não o
código efetivo. A segunda, mais forte e que constitui o subproduto metodológico
do trabalho, é que a LOC é sensível ao estilo de formatação e ao idioma
disponível em cada linguagem. A coluna pós-normalização da
Tabela~\ref{tab:prelim-loc} fornece a evidência direta desta segunda afirmação,
ao exibir o efeito isolado da formatação sobre a mesma base de código.

A análise das causas distingue-as por alcance. Duas operam \emph{entre}
linguagens e são eliminadas pela restrição intralinguagem. A primeira é o estilo
de formatação: no excerto do TCC~I, as vinculações que fixam a ordem de
avaliação apareciam em uma única linha, enquanto no código efetivo, formatado de
modo convencional, ocupam três linhas por operador binário. A segunda é o idioma
disponível: fixar a ordem de avaliação custa linhas em OCaml e é gratuito em
Haskell, onde o estilo aplicativo já sequencia à esquerda --- observe-se, porém,
que essa causa se cancela \emph{dentro} do OCaml, pois a referência imperativa
paga exatamente as mesmas linhas. A terceira causa é a única que opera dentro de
uma única linguagem: por manipular explicitamente a referência mutável e forçar
os resultados aritméticos, a linha de base do Haskell recorre à notação
\texttt{do} com vinculações intermediárias. É por isso que restringir a
comparação a pares da mesma linguagem elimina o confundimento com a sintaxe, mas
não o confundimento com o idioma escolhido dentro de cada linguagem.

Disso decorrem três ajustes na metodologia. Primeiro, a normalização por
formatadores, antes prevista como cuidado complementar, passa a ser condição de
validade de qualquer contagem, com perfil e versão fixados
(Seção~\ref{sec:metricas}). Segundo, as comparações de métricas de código passam
a ser restritas a pares da mesma linguagem, espelhando a normalização já adotada
para o desempenho. Terceiro, e porque a restrição intralinguagem não basta, o
peso da evidência sobre legibilidade desloca-se de LOC para a contagem de
operações de efeito e para a profundidade de aninhamento, que dependem menos do
\emph{estilo de formatação} --- embora permaneçam sensíveis ao idioma, limitação
que a restrição intralinguagem atenua sem eliminar. A hipótese H1 foi
reformulada nesses termos, e a definição de operação de efeito foi ampliada para
abranger também a realização direta, sem o que a métrica daria zero para as
linhas de base e a hipótese não seria confirmável. A LOC permanece reportada,
por transparência e comparabilidade com a literatura, mas deixa de sustentar
sozinha qualquer conclusão sobre QP1.

Registre-se que esse resultado ilustra, sobre um caso de uso concreto, que uma
métrica largamente empregada como \emph{proxy} de legibilidade pode inverter de
sinal por efeito de formatação e de idioma, mesmo quando a comparação é feita
dentro de uma única linguagem.

\subsection{Instrumentos de Coleta}\label{sec:instrumentos}

A Tabela~\ref{tab:metrics} consolidará as métricas de código, sempre em pares
intralinguagem, e a Tabela~\ref{tab:perf}, os resultados de desempenho como
razões de sobrecusto sobre a linha de base de cada linguagem.

\begin{table}[ht]
  \centering
  \caption{Instrumento de coleta das métricas de código, por par intralinguagem
           (valores a coletar após normalização).}
  \label{tab:metrics}
  \small\setlength{\tabcolsep}{4pt}
  \begin{tabular}{lcccc}
    \toprule
    & \multicolumn{2}{c}{\textbf{OCaml}} & \multicolumn{2}{c}{\textbf{Haskell}} \\
    \cmidrule(lr){2-3}\cmidrule(lr){4-5}
    \textbf{Métrica} & \textbf{Imperativa} & \textbf{Efeitos} & \textbf{\texttt{IORef}} & \textbf{Mônadas} \\
    \midrule
    LOC de \texttt{eval} (normalizadas)   & -- & -- & -- & -- \\
    Operações de efeito (total)           & -- & -- & -- & -- \\
    Operações de efeito (por construtor)  & -- & -- & -- & -- \\
    Profundidade de aninhamento           & -- & -- & -- & -- \\
    Boilerplate (LOC fora de \texttt{eval})  & -- & -- & -- & -- \\
    Instâncias de reutilização            & -- & -- & -- & -- \\
    Pontos de modificação (4º efeito)     & -- & -- & -- & -- \\
    Adições puras (4º efeito)             & -- & -- & -- & -- \\
    \bottomrule
  \end{tabular}
\end{table}

\begin{table}[ht]
  \centering
  \caption{Instrumento de coleta de desempenho: sobrecusto relativo à linha de
           base de cada linguagem, com intervalo de confiança por
           \textit{bootstrap} (valores a coletar).}
  \label{tab:perf}
  \small
  \begin{tabular}{lcccc}
    \toprule
    & \multicolumn{2}{c}{\textbf{Efeitos / imperativa (OCaml)}}
    & \multicolumn{2}{c}{\textbf{Mônadas / \texttt{IORef} (Haskell)}} \\
    \cmidrule(lr){2-3}\cmidrule(lr){4-5}
    \textbf{Cenário} & \textbf{Razão (IC 95\%)} & \textbf{Alocação}
                     & \textbf{Razão (IC 95\%)} & \textbf{Alocação} \\
    \midrule
    \textsc{Arith}  & -- & -- & -- & -- \\
    \textsc{State}  & -- & -- & -- & -- \\
    \textsc{Exn}    & -- & -- & -- & -- \\
    \textsc{Output} & -- & -- & -- & -- \\
    \bottomrule
  \end{tabular}
\end{table}

\subsection{Previsões}\label{sec:previsoes}

Três previsões orientam a análise. Primeiro, espera-se que o avaliador com
\textit{handlers} exiba contagem de operações de efeito e aninhamento próximos
aos da referência imperativa em OCaml, enquanto a versão monádica excederá os da
sua base em Haskell (H1). Segundo, espera-se que a inserção do quarto efeito
exija menos pontos de modificação na versão com efeitos, relativamente à sua
base, do que na versão monádica relativamente à dela (H2). Terceiro, espera-se
que o sobrecusto dos efeitos nativos sobre a base em OCaml seja inferior ao
sobrecusto da pilha de transformadores sobre a base em Haskell, sobretudo no
cenário \textsc{State} (P1); por comparar duas razões obtidas em
\textit{runtimes} distintos, essa previsão é exploratória.

Registre-se que a evidência já disponível tensiona parte dessas previsões. A
contagem de LOC (Tabela~\ref{tab:prelim-loc}) mostra a versão monádica mais
curta que a sua própria linha de base, e a Tabela~\ref{tab:artefatos} mostra, por
arquivo, a versão com \textit{handlers} substancialmente maior que a base em
OCaml, o que aponta contra H3 na métrica de \textit{boilerplate}. Ambos os
sinais serão confrontados com as métricas normalizadas na fase seguinte.

%% ============================================================
\section{Atividades Desenvolvidas e Cronograma}\label{sec:atividades}
%% ============================================================

Desde o TCC~I \cite{marchioni:tcc1}, consolidaram-se a formalização da semântica
e do critério de equivalência (Seção~\ref{sec:semantica}), as quatro
implementações e o gerador de cargas (Seção~\ref{sec:impl}), a suíte de oráculo
(Seção~\ref{sec:oraculo}), os dois \textit{harnesses} (Seção~\ref{sec:harness})
e o ferramental de medição (Seção~\ref{sec:organizacao}). Duas divergências do
relatório anterior foram corrigidas antes da implementação --- o uso de
\texttt{div} em lugar de \texttt{quot} e a ordem das camadas da pilha monádica,
que descartava o estado no erro --- e a própria implementação expôs dois
problemas que a especificação em papel não revelava: a ordem de avaliação não
especificada dos argumentos em OCaml e o tratamento ausente de variável não
vinculada na versão com efeitos. Todos estão hoje cobertos por casos de teste.
A revisão de literatura foi atualizada com
\cite{schrijvers:19,xie:20,karachalias:21,gaissert:25}, e a contagem de LOC
sobre o código efetivo motivou a restrição das comparações a pares
intralinguagem (Seção~\ref{sec:refinamento}).

As atividades restantes são: \textbf{A1}, validação no ambiente definitivo,
resolvendo as pendências da Seção~\ref{sec:estado}; \textbf{A2}, publicação do
repositório com depósito arquivado e identificador persistente; \textbf{A3},
inserção do efeito \texttt{Input} nas quatro versões, em ramo separado, com
registro dos pontos de modificação para responder a QP2; \textbf{A4}, coleta
das métricas de código por dupla verificação independente do autor e do
orientador; \textbf{A5}, execução dos quatro cenários nas condições da
Seção~\ref{sec:protocolo}, com 30 amostras por configuração, três sementes e
intervalos de confiança por \textit{bootstrap}; \textbf{A6}, redação de
Resultados e Discussão, preenchendo as Tabelas~\ref{tab:metrics}
e~\ref{tab:perf} e interpretando-as à luz de H1--H4 e P1; e \textbf{A7},
revisão final e reexecução em ambiente limpo para confirmar a
\emph{repetibilidade} --- a \emph{reprodutibilidade} por terceiros é viabilizada
pelos artefatos publicados em A2. A Tabela~\ref{tab:cronograma} apresenta o
cronograma e as dependências entre elas. O principal risco é o \textit{switch}
com \textit{flambda} não ser construído com êxito, cuja resposta está
registrada na Seção~\ref{sec:estado}.

\begin{table}[ht]
  \centering
  \caption{Cronograma das próximas atividades, com dependências. Uma atividade
           só se inicia após a conclusão daquelas de que depende. A entrega para
           a banca ocorre em 23/11 e a defesa entre 02 e 10/12.}
  \label{tab:cronograma}
  \small\setlength{\tabcolsep}{4pt}
  \begin{tabular}{llcccc}
    \toprule
    \textbf{Atividade} & \textbf{Dep.} & \textbf{Set} & \textbf{Out} & \textbf{Nov} & \textbf{Dez} \\
    \midrule
    A1. Validação no ambiente definitivo   & ---     & $\bullet$ &           &           &           \\
    A2. Publicação do repositório          & A1      & $\bullet$ &           &           &           \\
    A3. Implementação do quarto efeito     & A1      & $\bullet$ & $\bullet$ &           &           \\
    A4. Coleta de métricas de código       & A3      &           & $\bullet$ &           &           \\
    A5. Execução dos \textit{benchmarks}   & A1      &           & $\bullet$ &           &           \\
    A6. Redação de Resultados e Discussão  & A4,\,A5 &           & $\bullet$ & $\bullet$ &           \\
    A7. Revisão final e entrega à banca    & A6      &           &           & $\bullet$ &           \\
    Defesa                                 & A7      &           &           &           & $\bullet$ \\
    \bottomrule
  \end{tabular}
\end{table}

%% ============================================================
\section{Considerações Parciais}\label{sec:consideracoes}
%% ============================================================

Este artigo reportou o estado do trabalho no primeiro ponto de controle do
TCC~II. Em relação ao relatório do TCC~I \cite{marchioni:tcc1}, três avanços se
destacam.

\paragraph{Passagem do desenho à implementação.} As quatro versões do
interpretador, a suíte de oráculo e os \textit{harnesses} de \textit{benchmark}
existem como código. Duas divergências foram identificadas pela releitura
crítica do relatório anterior --- a operação de divisão inteira e a ordem das
camadas monádicas ---, e duas outras só se revelaram no ato de implementar, que
a especificação em papel não expunha: a ordem de avaliação não especificada dos
argumentos em OCaml e o tratamento de variáveis não vinculadas na versão com
efeitos. Todas as quatro foram corrigidas e estão hoje cobertas por casos de
teste que falham se a decisão for revertida.

\paragraph{Construção de instrumentos auditáveis.} Um gerador de cargas
replicado nas duas linguagens, com teste automático de conformidade entre as
réplicas, assegura que Haskell e OCaml avaliem as mesmas árvores de sintaxe; e
uma suíte de oráculo derivada da semântica formal condiciona a execução dos
\textit{benchmarks} à equivalência funcional verificada.

\paragraph{Refinamento metodológico motivado por evidência.} A contagem de LOC
sobre o código efetivo inverteu a tendência observada nos excertos do TCC~I. A
análise das causas mostrou que a métrica é sensível ao estilo de formatação, ao
custo sintático de fixar a ordem de avaliação em OCaml e ao idioma com que cada
linha de base manipula o estado --- nenhum dos quais diz respeito ao mecanismo de
efeitos. Em resposta, as comparações de código passaram a ser restritas a pares
da mesma linguagem, pela mesma razão que já levava o desempenho a ser normalizado
contra uma linha de base por linguagem.

\paragraph{Limitações e etapas seguintes.} As conclusões que o trabalho poderá
oferecer restringem-se a um único caso de uso, a versões específicas de
compilador e hardware, ao regime de abortagem total do efeito de exceção, e às
quatro dimensões adotadas --- ficando de fora a segurança estática de efeitos,
em que as abordagens diferem substancialmente. A comparação direta entre as duas
linguagens permanece exploratória, por ser confundida pelo \textit{runtime}.
Registre-se ainda que a evidência preliminar sobre o código já tensiona parte
das hipóteses: em Haskell, a versão monádica é mais curta que a sua própria
linha de base, e em OCaml o custo estrutural dos \textit{handlers}, medido por
arquivo, supera o da referência imperativa. As medições que permitirão responder
em definitivo às quatro questões de pesquisa constituem a etapa seguinte, e a
validação das implementações no ambiente definitivo é o pré-requisito imediato.

%% ============================================================
%% CONFERIR antes da entrega: a lista de secoes abaixo e a versao do modelo
%% devem refletir exatamente o uso que voce fez da ferramenta (Arts. 6o a 8o).
\section*{Agradecimentos}

Em conformidade com os Arts.~6\textsuperscript{o} a 8\textsuperscript{o} do
Regulamento Geral de TCC do curso, declara-se o uso de ferramenta de
Inteligência Artificial neste trabalho. Foi utilizado o modelo Claude
(Anthropic), por meio da interface Claude Code, com a finalidade de revisão
crítica do texto: verificação de consistência entre seções, conferência da
correspondência entre afirmações e referências citadas, identificação de
redundâncias e sugestões de redação e de concisão. A ferramenta contribuiu nas
Seções~\ref{sec:intro}, \ref{sec:impl} e~\ref{sec:atividades}. A concepção do
trabalho, o desenho experimental, a formalização da semântica, as
implementações e a interpretação dos resultados são de autoria do aluno, que
permanece integralmente responsável pelo conteúdo apresentado.

%% ============================================================
\bibliographystyle{sbc}
\bibliography{sbc-template}

\end{document}
